\documentclass[11pt]{amsart}

\usepackage[utf8]{inputenc}

\usepackage{amsmath,amssymb,amsthm,mathtools}
\usepackage[margin=1.05in]{geometry}
\usepackage{enumitem}
\usepackage{microtype}
\usepackage{xcolor}
\usepackage{listings}
\usepackage[colorlinks=true,linkcolor=blue!50!black,citecolor=blue!50!black,urlcolor=blue!50!black]{hyperref}

\lstset{basicstyle=\ttfamily\scriptsize,breaklines=true,columns=fullflexible,
  frame=single,framerule=0.3pt,rulecolor=\color{black!40},showstringspaces=false,
  keepspaces=true}

\newtheorem{theorem}{Theorem}
\renewcommand{\thetheorem}{\Alph{theorem}}
\newtheorem{maincor}[theorem]{Corollary}
\newtheorem{lemma}{Lemma}[section]
\newtheorem{proposition}[lemma]{Proposition}
\newtheorem{corollary}[lemma]{Corollary}
\theoremstyle{definition}
\newtheorem{definition}[lemma]{Definition}
\newtheorem{problem}[lemma]{Problem}
\newtheorem{principle}[lemma]{Principle}
\newtheorem{question}[lemma]{Question}
\theoremstyle{remark}
\newtheorem{remark}[lemma]{Remark}
\newtheorem{example}[lemma]{Example}

\newcommand{\N}{\mathbb{N}}
\newcommand{\Z}{\mathbb{Z}}
\newcommand{\Q}{\mathbb{Q}}
\newcommand{\R}{\mathbb{R}}
\DeclareMathOperator{\lcm}{lcm}

\title[An exponent window for Erd\H{o}s Problem \#265]{An Exponent Window for
Erd\H{o}s Problem \#265:\\ Cross-Tail Integrality and Controlled\\ Two-Target Greedy Expansions}

\author{Keio}
\thanks{Research draft, October 28, 2026. The proofs and the exposition were developed
interactively with an AI assistant (Claude, Anthropic); the author is responsible for
the content. Every algebraic identity, integrality statement and the construction of
Theorem~\ref{thm:C} are accompanied by an exact-arithmetic verification script
(Appendix~\ref{app:scripts}). \emph{Priority:} the special case of
Theorem~\ref{thm:A} that answers the limsup question of Erd\H{o}s Problem \#265 was
established first by K.~Kitamura in a Lean~4 formalization \cite{Kit26}; see
Section~\ref{sec:lit}. This draft has not been peer reviewed.}

\date{October 28, 2026}

\begin{document}

\begin{abstract}
Let $a_1<a_2<\cdots$ be integers. Erd\H{o}s and Graham asked how fast $(a_n)$ can grow
if both $\sum 1/a_n$ and $\sum 1/(a_n-1)$ are rational (Erd\H{o}s Problem \#265).
Kova\v{c} and Tao proved that doubly exponential growth is possible (their method gives
$a_n^{1/\beta^n}\to\infty$ for every $\beta<\sqrt{6/5}=1.0954\ldots$), and K.~Kitamura
recently gave a Lean-verified proof of Erd\H{o}s's prediction $a_n^{1/2^n}\to1$. We study the growth of
sequences with two or more rational translated reciprocal sums.

(1) We give an independent, human-readable proof that $a_n^{1/2^n}\to1$ whenever two
distinct \emph{rational} translates (or two non-proportional rational affine images)
of $(a_n)$ have rational reciprocal sums. The key input is the integrality of the
\emph{cross tail} $E_n=R_nR_n'-(R_n-R_n')/\Delta=\sum_{j\ne k\ge n}1/(\alpha_j\beta_k)$;
for Erd\H{o}s's pair it coincides, up to the unit factor $1-R'_n$, with the ``second
residual'' of Kitamura's formalization. We combine it with a frozen-jump lemma and a
checkpoint argument.

(2) A multilinear version, based on M\"obius inversion on the partition lattice, shows
that $d$ rational integer translates force $\log a_n=O((d+1)^{n/d})$ whenever the tails
are dominated by their first terms (e.g.\ if $\liminf a_{n+1}/a_n>1$). For Problem \#265
this gives $\log a_n=O(3^{n/2})$ for such sequences; this case can also be derived in a
few lines from lemmas in Kitamura's formalization, while $d\ge3$ needs the multilinear
integrality.

(3) Conversely, a \emph{controlled two-target greedy} construction produces, for every
$1<\beta<\beta_*:=(\sqrt{13}-1)/2=1.3027\ldots$, sequences with both sums rational,
$a_{n+1}/a_n\to\infty$ and $\log a_n\sim c\,\beta^n$. This improves the exponent
$\sqrt{6/5}$ of Kova\v{c} and Tao.

Hence the optimal growth exponent for Problem \#265 among sequences with
$\liminf a_{n+1}/a_n>1$ lies in $[\beta_*,\sqrt3\,]=[1.3027\ldots,1.7320\ldots]$; a
counting heuristic suggests that $\beta_*$ is the true value.
\end{abstract}

\maketitle

\tableofcontents

%======================================================================
\section{Introduction}\label{sec:intro}
%======================================================================

\subsection{Erd\H{o}s Problem \#265}
An \emph{Ahmes series} is a convergent series $\sum 1/a_n$ of reciprocals of a strictly
increasing sequence of positive integers. A classical folklore criterion states that
$\sum 1/a_n$ is irrational as soon as $a_n^{1/2^n}\to\infty$, and Sylvester's sequence
\begin{equation}\label{eq:sylvester}
s_0=2,\qquad s_{n+1}=s_n^2-s_n+1,\qquad \sum_{n\ge0}\frac1{s_n}=1,
\end{equation}
shows that this is sharp: here $s_n^{1/2^n}\to E^2=1.5979\ldots$ with
Vardi's constant $E=1.2640847\ldots$ \cite{Sylvester1880,Vardi91}.

Erd\H{o}s and Graham \cite[p.~64]{ErGr80} (see also \cite[pp.~104--105]{Er88} and
\cite[p.~334]{Er85}) asked how fast $(a_n)$ can grow if the two series
\[
\sum_n\frac1{a_n}\qquad\text{and}\qquad \sum_n\frac1{a_n-1}
\]
are \emph{both} rational. David Cantor observed that $a_n=\binom n2$ $(n\ge3)$ works,
with sums $1$ and $11/9$. Erd\H{o}s wrote that he was sure that $a_n^{1/n}\to\infty$ is
possible but that $a_n^{1/2^n}$ must tend to~$1$ \cite[pp.~104--105]{Er88}. This is
Problem~\#265 of Bloom's database \cite{Bloom265}.

Kova\v{c} \cite{Kov24} constructed exponentially growing examples, and Kova\v{c} and
Tao \cite{KT25} proved that for every $d$ there are sequences with
$\sum 1/(a_n+j)\in\Q$ for $0\le j<d$ and $a_n^{1/\beta^n}\to\infty$ for some $\beta>1$;
their proof works for every $\beta<((2d+2)/(2d+1))^{1/d}$ \cite[(7.9)]{KT25}, which for
two series is $\beta<\sqrt{6/5}$. This confirms
the first half of Erd\H{o}s's prediction in a strong form. The second half---equivalently,
that $\limsup a_n^{1/2^n}>1$ is impossible---was proved recently by K.~Kitamura in a
Lean~4 formalization \cite{Kit26}, which was independently checked in the course of a
review for the \emph{formal-conjectures} project \cite{FC6839}; at the time of writing
the database still lists the problem as open, the remaining question being the precise
exponent \cite{Bloom265}. Kova\v{c} and Tao also remark that ``it remains unclear what
the optimal exponent of the doubly exponential growth is'' \cite{KT25}.

This paper is about that remaining question, and about its analogue for several
translates. For sequences whose tails are dominated by their first terms we show that
the optimal exponent lies between $\beta_*=1.3027\ldots$ and $\sqrt3$
(Corollary~\ref{cor:D}). The lower bound is new and improves the Kova\v{c}--Tao exponent
$\sqrt{6/5}=1.0954\ldots$. The upper bound is essentially implicit in the lemmas of
\cite{Kit26}, and we extend it to $d$ translates. We also give an independent proof of the endpoint
$a_n^{1/2^n}\to1$ which works for two arbitrary rational translates.

\subsection{A problem: growth under rational translates}

\begin{problem}[Growth under rational translates]\label{prob:shifts}
Let $a_1<a_2<\cdots$ be integers and let $t_1,\dots,t_d$ be distinct rational numbers
with $a_n+t_i\ne0$ for all $n,i$. Suppose that
\[
\sum_n\frac{1}{a_n+t_i}\in\Q\qquad(i=1,\dots,d).
\]
How fast can $(a_n)$ grow? In particular, determine
\[
\beta_d^{\mathrm{reg}}:=\sup\Bigl\{\beta>1:\ \exists\,(a_n)\ \text{as above with}\
\liminf_{n\to\infty}\frac{a_{n+1}}{a_n}>1\ \text{and}\ a_n^{1/\beta^n}\to\infty\Bigr\}
\]
for the consecutive shifts $t_i=i-1$.
\end{problem}

For $d=1$ the answer is classical: $c^{2^n}$-growth with $c>1$ is possible (Sylvester)
and $a_n^{1/2^n}\to\infty$ is impossible, so $\beta_1^{\mathrm{reg}}=2$. Erd\H{o}s's
problem is the case $d=2$ (after the substitution $a_n\mapsto a_n-1$).

\subsection{Main results}

\begin{theorem}[Two rational translates]\label{thm:A}
Let $a_1<a_2<\cdots$ be integers and let $t\ne t'$ be rational numbers with
$a_n+t\ne 0\ne a_n+t'$ for all $n$. If $\sum_n 1/(a_n+t)$ and $\sum_n 1/(a_n+t')$ are
both rational, then
\[
\lim_{n\to\infty}\frac{\log a_n}{2^n}=0,\qquad\text{equivalently}\qquad
\lim_{n\to\infty}a_n^{1/2^n}=1 .
\]
\end{theorem}

For $\{t,t'\}=\{0,-1\}$ this is Kitamura's theorem \cite{Kit26}; see
Section~\ref{sec:lit} for a comparison of the two proofs.

\begin{theorem}[Regular growth under $d$ rational translates]\label{thm:B}
Let $d\ge1$, let $h_1<\cdots<h_d$ be integers, and let $a_1<a_2<\cdots$ be positive
integers with $a_n+h_1\ge1$ and $\sum_n 1/(a_n+h_i)\in\Q$ for $i=1,\dots,d$. Assume that the
tails are dominated by their first terms,
\begin{equation}\label{eq:tails}
\sup_{n\ge1}\ a_n\sum_{k\ge n}\frac1{a_k}<\infty ,
\end{equation}
which holds in particular under the regularity condition
\begin{equation}\label{eq:regular}
\liminf_{n\to\infty}\frac{a_{n+1}}{a_n}>1 .
\end{equation}
Then there is a constant $C$ such that
\[
\sum_{j=0}^{d-1}\log a_{n+j}\le d\sum_{k<n}\log a_k+C\quad(n\ge1),
\qquad\text{hence}\qquad \log a_n=O\bigl((d+1)^{n/d}\bigr).
\]
\end{theorem}

For $d=1$ this is the folklore bound, attained by Sylvester's sequence; for $d=2$ it
gives $\log a_n=O(3^{n/2})$. For Erd\H{o}s's pair the case $d=2$ is not stated in
\cite{Kit26}, but it follows in a few lines from three lemmas proved there (positivity of
the second residual, the bound $\le8T_nT_{n+1}$, and its scaled integrality), so we do not
claim it as new; the case $d\ge3$ requires the multilinear integrality of
Section~\ref{sec:B}. By the scaling of Section~\ref{subsec:rational}, which preserves
\eqref{eq:regular}, the same holds for rational shifts.

\begin{theorem}[A construction]\label{thm:C}
Let $1<\beta<\beta_*:=\frac{\sqrt{13}-1}{2}=1.30277\ldots$, the positive root of
$\beta^2+\beta=3$. There is a strictly increasing sequence of integers $a_n\ge2$ such
that $\sum_n1/a_n$ and $\sum_n1/(a_n-1)$ are both rational, $a_{n+1}/a_n\to\infty$, and
\[
\lim_{n\to\infty}\frac{\log a_n}{\beta^n}=c\qquad\text{for some }c\in(0,\infty).
\]
\end{theorem}

\begin{maincor}[An exponent window for Problem \#265]\label{cor:D}
Let $2\le a_1<a_2<\cdots$ be integers such that $\sum1/a_n$ and $\sum1/(a_n-1)$ are both
rational. Then $a_n^{1/2^n}\to 1$ \textup{(}\cite{Kit26}; Theorem~\ref{thm:A}\textup{)}.
If moreover $\liminf a_{n+1}/a_n>1$, then $\log a_n=O(3^{n/2})$ \textup{(}Theorem~\ref{thm:B};
cf.\ \cite{Kit26}\textup{)}. Conversely, for every
$\beta<\beta_*$ there are such sequences with $a_{n+1}/a_n\to\infty$ and
$a_n^{1/\beta^n}\to\infty$. In the notation of Problem~\ref{prob:shifts},
\[
1.3027\ldots=\frac{\sqrt{13}-1}{2}\ \le\ \beta^{\mathrm{reg}}_2\ \le\ \sqrt3=1.7320\ldots .
\]
\end{maincor}

\begin{maincor}[At most one rational translate]\label{cor:E}
Let $a_1<a_2<\cdots$ be integers with $\limsup a_n^{1/2^n}>1$. Then there is at most
one rational number $t$ with $a_n+t\ne0$ for all $n$ and $\sum_n 1/(a_n+t)\in\Q$.
\end{maincor}

For sequences satisfying an exact quadratic recursion such as Sylvester's, this also
follows from Duverney's criterion \cite{Duverney01}; see Example~\ref{ex:sylvester}.

\subsection{The methods}
Classical irrationality arguments for Ahmes series are \emph{first order}: if
$\sum 1/a_n=u/q$, then $q\,\lcm(a_1,\dots,a_{n-1})\cdot\sum_{k\ge n}1/a_k$ is a positive
integer, and when the tail is comparable to $1/a_n$ this gives
$a_n\lesssim a_1\cdots a_{n-1}$, i.e.\ growth at most $c^{2^n}$. Two rational
translates give three integral quantities at each level $n$: the two scaled tails and
the scaled tail of the \emph{defect series} $\sum 1/(\alpha_k\beta_k)$. A fourth,
\emph{bilinear}, combination is also integral:

\begin{principle}[Cross-tail integrality]\label{princ}
Let $R_n=\sum_{k\ge n}1/\alpha_k$ and $R'_n=\sum_{k\ge n}1/\beta_k$ with
$\beta_k=\alpha_k+\Delta$. The off-diagonal part of the product of the two tails,
\[
E_n=R_nR_n'-\frac{R_n-R_n'}{\Delta}=\sum_{\substack{j,k\ge n\\ j\ne k}}\frac{1}{\alpha_j\beta_k},
\]
is a positive rational number whose denominator divides
$\Delta Q^2\prod_{k<n}\alpha_k\beta_k$ (Lemma~\ref{lem:cross}). Its size is at most
$2R_nR_{n+1}$, i.e.\ it is governed by two \emph{consecutive} tails.
\end{principle}

For $\Delta=1$, $\Delta E_n=(1-\Delta R'_n)\bigl(R'_n/(1-\Delta R'_n)-R_n\bigr)$, and the
second factor is exactly the ``second residual'' $T_n/(1-T_n)-V_n$ used in Kitamura's
formal proof \cite{Kit26} (Remark~\ref{rem:residual}); so for Erd\H{o}s's pair the
cross-tail integrality is equivalent to Kitamura's second-residual integrality. What is
new here is the use of $E_n$ for arbitrary $\Delta$ and for affine images, its
multilinear extension, and the quantitative consequences below.

When the tails are dominated by their first terms, Principle~\ref{princ} yields
$a_na_{n+1}\lesssim(a_1\cdots a_{n-1})^2$, a second-order recursion with growth rate
$\sqrt3<2$; this is the mechanism behind Theorem~\ref{thm:B} for $d=2$ (the same
inequality can be extracted from the lemmas of \cite{Kit26}). In general the
tails need not be dominated by their first terms. Such irregularities are controlled by a
\emph{frozen-jump lemma} (Lemma~\ref{lem:frozen}), obtained from the integrality of the
defect series, and a checkpoint argument for the normalized quantity $\pi_n/2^n$ (where
$\pi_n$ is the logarithm of the partial product) completes the proof of
Theorem~\ref{thm:A}. The multilinear version of Principle~\ref{princ}
(Lemma~\ref{lem:multi}) uses M\"obius inversion on the lattice of set partitions and
partial fractions.

The construction of Theorem~\ref{thm:C} is a \emph{controlled two-target greedy
algorithm}. We prescribe the two rational sums in advance and choose the terms one by
one. Integrality is then automatic, and the only issue is to keep the two remaining
tails compatible. We measure compatibility by $\Gamma_n=\tau_n^2-\delta_n$, where
$\tau_n=\sum_{k\ge n}1/a_k$ and $\delta_n=\sum_{k\ge n}1/(a_k(a_k-1))$. A single integer
choice per step can steer $\Gamma_{n+1}$ to any prescribed size down to the scale
$a_n^{-3}$. For this greedy one therefore needs $a_{n+1}a_{n+2}\lesssim a_n^3$, i.e.\
$\beta^2+\beta\le3$, and the algorithm achieves every $1<\beta<\beta_*$
(Lemma~\ref{lem:step}).

\subsection{Organization and verification}
Section~\ref{sec:lemmas} proves the integrality lemmas and their logarithmic
consequences, Section~\ref{sec:A} proves Theorem~\ref{thm:A}, Section~\ref{sec:B}
proves Theorem~\ref{thm:B}, and Section~\ref{sec:C} proves Theorem~\ref{thm:C}.
Section~\ref{sec:cor} derives the corollaries and discusses related Erd\H{o}s problems,
Section~\ref{sec:limits} discusses the limits of the method and open problems, and
Section~\ref{sec:lit} compares with the literature. Appendix~\ref{app:scripts} contains
a script that checks every algebraic identity and integrality statement of
Sections~\ref{sec:lemmas} and~\ref{sec:B} in exact rational arithmetic, and runs the
construction of Section~\ref{sec:C} for several steps while checking every conclusion of
Lemma~\ref{lem:step}.

%======================================================================
\section{Tails, partial products and integrality}\label{sec:lemmas}
%======================================================================

\subsection{Setting}\label{subsec:setting}
Throughout Sections~\ref{sec:lemmas} and~\ref{sec:A} we fix integers $h<h'$ and a
strictly increasing sequence of integers $(a_n)_{n\ge1}$ with $a_n+h\ge1$ for all $n$,
such that
\[
S:=\sum_{n\ge1}\frac{1}{a_n+h}\in\Q,\qquad S':=\sum_{n\ge1}\frac1{a_n+h'}\in\Q .
\]
We write
\[
\alpha_k:=a_k+h,\qquad \beta_k:=a_k+h'=\alpha_k+\Delta,\qquad \Delta:=h'-h\ge1 .
\]

\emph{Normalization.} Removing finitely many initial terms does not affect the
hypotheses (the new sums are still rational) and does not affect the conclusion of
Theorem~\ref{thm:A}, since $\log a_{n+N}=o(2^{n})$ is equivalent to $\log a_n=o(2^n)$.
We may therefore assume that
\begin{equation}\label{eq:normal}
S=\sum_{k\ge1}\frac1{\alpha_k}<1,\qquad\text{so that}\quad \alpha_k\ge2\ \text{ for all }k.
\end{equation}
After this normalization we fix an integer $Q\ge1$ with $QS,QS'\in\Z$.

For $n\ge1$ we define the partial products and tails
\[
\Pi_n:=\prod_{k<n}\alpha_k,\qquad \Pi'_n:=\prod_{k<n}\beta_k,\qquad
R_n:=\sum_{k\ge n}\frac1{\alpha_k},\qquad R'_n:=\sum_{k\ge n}\frac1{\beta_k},
\]
the \emph{defect tails} and the \emph{cross tails}
\[
D_n:=\sum_{k\ge n}\frac{1}{\alpha_k\beta_k},\qquad
E_n:=\sum_{\substack{j,k\ge n\\ j\neq k}}\frac1{\alpha_j\beta_k}.
\]
All series converge absolutely, and $0<R'_n<R_n\le R_1<1$.

\begin{lemma}[Algebraic identities]\label{lem:identities}
For every $n\ge1$,
\[
D_n=\frac{R_n-R'_n}{\Delta},\qquad E_n=R_nR'_n-D_n .
\]
\end{lemma}

\begin{proof}
Since $\beta_k-\alpha_k=\Delta$, we have $\frac{1}{\alpha_k\beta_k}=
\frac1\Delta\bigl(\frac1{\alpha_k}-\frac1{\beta_k}\bigr)$; summing over $k\ge n$ gives the
first identity. Expanding the absolutely convergent product
$R_nR'_n=\sum_{j,k\ge n}1/(\alpha_j\beta_k)$ and separating the diagonal $j=k$ gives
$R_nR'_n=D_n+E_n$.
\end{proof}

\subsection{Three integrality lemmas}

\begin{lemma}[Single tails]\label{lem:single}
For every $n\ge1$ the numbers $X_n:=Q\Pi_nR_n$ and $X'_n:=Q\Pi'_nR'_n$ are positive
integers.
\end{lemma}

\begin{proof}
$X_n=\Pi_n\cdot QS-\sum_{k<n}Q\,\Pi_n/\alpha_k$, and each $\Pi_n/\alpha_k$ $(k<n)$ is an
integer; $X_n>0$ because $R_n>0$. The same argument applies to $X'_n$.
\end{proof}

\begin{lemma}[Defect integrality]\label{lem:defect}
For every $n\ge1$,
\[
N_n:=\Delta Q\,\Pi_n\Pi'_n\,D_n=\Pi'_nX_n-\Pi_nX'_n
\]
is a positive integer.
\end{lemma}

\begin{proof}
By Lemma~\ref{lem:identities},
$\Delta Q\Pi_n\Pi'_nD_n=Q\Pi_n\Pi'_n(R_n-R'_n)=\Pi'_nX_n-\Pi_nX'_n\in\Z$
(Lemma~\ref{lem:single}), and $D_n>0$.
\end{proof}

\begin{lemma}[Cross-tail integrality]\label{lem:cross}
For every $n\ge1$,
\[
M_n:=\Delta Q^2\,\Pi_n\Pi'_n\,E_n=\Delta X_nX'_n-QN_n
\]
is a positive integer.
\end{lemma}

\begin{proof}
By Lemma~\ref{lem:identities},
$\Delta Q^2\Pi_n\Pi'_nE_n=\Delta(Q\Pi_nR_n)(Q\Pi'_nR'_n)-Q\cdot\Delta Q\Pi_n\Pi'_nD_n
=\Delta X_nX'_n-QN_n$, which is an integer by Lemmas~\ref{lem:single}
and~\ref{lem:defect}. It is positive because the off-diagonal sum $E_n$ contains the
positive term $1/(\alpha_n\beta_{n+1})$ (the sequence is infinite).
\end{proof}

\begin{remark}
Lemma~\ref{lem:cross} is the cross-tail integrality principle
(Principle~\ref{princ}). Note that $M_n\ge1$ is \emph{not} a consequence of the sizes of
$X_n,X'_n,N_n$: these three integers can be large while $M_n$ is forced to be at least
$1$ by a cancellation between $\Delta X_nX'_n$ and $QN_n$. This cancellation is what
couples two consecutive scales.
\end{remark}

\begin{remark}[Residual form; relation to \cite{Kit26}]\label{rem:residual}
Whenever $\Delta R'_n<1$ one has the identity
\[
\Delta E_n=(1-\Delta R'_n)\,\rho_n,\qquad \rho_n:=\frac{R'_n}{1-\Delta R'_n}-R_n ,
\]
since $(1-\Delta R'_n)\rho_n=R'_n-R_n+\Delta R_nR'_n$. If the tail consisted of a single
term $1/\beta_n$, then $\rho_n=1/(\beta_n-\Delta)-1/\alpha_n=0$; thus $\rho_n$ measures
how far the tail is from a single term. For $\Delta=1$ and the pair
$(\alpha_k,\beta_k)=(a_k-1,a_k)$ of Problem~\#265 we have $R'_n=T_n:=\sum_{k\ge n}1/a_k$
and $R_n=V_n:=\sum_{k\ge n}1/(a_k-1)$, and $\rho_n=T_n/(1-T_n)-V_n$ is the ``second
residual'' of \cite{Kit26}, whose scaled integrality is proved there. Thus
Lemma~\ref{lem:cross} generalizes the key integrality of \cite{Kit26} to arbitrary
$\Delta$.
\end{remark}

\subsection{Size bounds and logarithmic form}

\begin{lemma}[Size bounds]\label{lem:size}
For all $n\ge1$ and all $1\le m\le p$:
\begin{enumerate}[label=\textup{(\roman*)}]
\item $E_n<2R_nR_{n+1}$;
\item $D_m<\dfrac{p-m}{\alpha_m^2}+\dfrac{R_p}{\alpha_p}$;
\item $\Pi'_n<e^{\Delta}\,\Pi_n$.
\end{enumerate}
\end{lemma}

\begin{proof}
(i) Since $1/\beta_k<1/\alpha_k$,
\[
E_n<\sum_{\substack{j,k\ge n\\ j\neq k}}\frac1{\alpha_j\alpha_k}
=2\sum_{n\le j<k}\frac1{\alpha_j\alpha_k}=2\sum_{j\ge n}\frac{R_{j+1}}{\alpha_j}
\le 2R_{n+1}\sum_{j\ge n}\frac{1}{\alpha_j}=2R_nR_{n+1}.
\]
(ii) We have $D_m<\sum_{k\ge m}\alpha_k^{-2}$. For $m\le k<p$ we use
$\alpha_k\ge\alpha_m$, and for $k\ge p$ we use $\alpha_k^{-2}\le\alpha_p^{-1}\alpha_k^{-1}$.
(iii) $\Pi'_n/\Pi_n=\prod_{k<n}(1+\Delta/\alpha_k)\le\exp(\Delta\sum_k1/\alpha_k)<e^{\Delta}$
by \eqref{eq:normal}.
\end{proof}

We pass to logarithms. Put
\[
\ell_k:=\log\alpha_k>0,\qquad \pi_n:=\sum_{k<n}\ell_k=\log\Pi_n,\qquad t_n:=-\log R_n>0,
\]
and fix the constants
\[
\kappa_1:=\log\bigl(2\Delta Q^2e^{\Delta}\bigr),\qquad \kappa_2:=\log\bigl(2\Delta Qe^{\Delta}\bigr)>0 .
\]
Note that $\ell_k$ is strictly increasing, $t_n$ is non-decreasing and $t_n\to\infty$.

\begin{corollary}[Logarithmic inequalities]\label{cor:log}
For all $n\ge 1$ and all $1\le m\le p$ the following hold.
\begin{enumerate}[label=\textup{(\Alph*)}]
\item[\textup{(K)}] \emph{(cross-tail inequality)} $\ t_n+t_{n+1}<2\pi_n+\kappa_1$.
\item[\textup{(F)}] \emph{(frozen-jump dichotomy)} Either
$2\ell_m<2\pi_m+\log(p-m)+\kappa_2$, or $\ell_p+t_p<2\pi_m+\kappa_2$.
\textup{(}For $p=m$ the first alternative is void, $\log 0=-\infty$.\textup{)}
\item[\textup{(J)}] \emph{(jump bound)} $\ \ell_m<2\pi_m+\kappa_2$.
\end{enumerate}
\end{corollary}

\begin{proof}
(K) By Lemmas~\ref{lem:cross} and~\ref{lem:size}(i),(iii),
$1\le M_n<\Delta Q^2\,e^{\Delta}\,\Pi_n^2\cdot2R_nR_{n+1}$; take logarithms.

(F) By Lemma~\ref{lem:defect} and Lemma~\ref{lem:size}(iii),
$D_m\ge 1/(\Delta Q\Pi_m\Pi'_m)>e^{-2\pi_m}/(\Delta Qe^{\Delta})=2e^{-2\pi_m-\kappa_2}$.
Combining with Lemma~\ref{lem:size}(ii),
\[
2e^{-2\pi_m-\kappa_2}<(p-m)e^{-2\ell_m}+e^{-\ell_p-t_p},
\]
so one of the two summands exceeds $e^{-2\pi_m-\kappa_2}$, which is the stated
dichotomy. For $p=m$ the first summand vanishes.

(J) Take $p=m$ in (F) and use $t_m>0$.
\end{proof}

\begin{remark}
Inequality (J) is the defect-series analogue of the folklore bound: it allows a
single term to be as large as the square of the product of all previous terms. The
dichotomy (F) says that such a large term at index $m$ \emph{freezes} the following
terms: as long as $p-m$ is small compared with $e^{2(\ell_m-\pi_m)}$, all later terms
up to index $p$ stay below $e^{2\pi_m+\kappa_2}$.
\end{remark}

%======================================================================
\section{Proof of Theorem~\ref{thm:A}}\label{sec:A}
%======================================================================

We keep the setting of Section~\ref{subsec:setting}. The reduction of rational shifts
to integer shifts is given at the end of the section.

\subsection{Big indices and frozen windows}
Put
\[
\kappa:=\kappa_2+\log 3,\qquad v_n:=\frac{\pi_n}{2^n}\quad(n\ge1).
\]

\begin{definition}
An index $m$ is \emph{big} if $\ell_m>\pi_m+\kappa$.
\end{definition}

\begin{lemma}[Frozen-jump lemma]\label{lem:frozen}
Let $m$ be big. Then
\begin{enumerate}[label=\textup{(\alph*)}]
\item $\ell_p<2\pi_m+\kappa_2$ for $p=m,m+1,m+2,m+3$;
\item none of the indices $m+1,m+2,m+3$ is big.
\end{enumerate}
\end{lemma}

\begin{proof}
(a) For $p=m$ this is (J). Let $p\in\{m+1,m+2,m+3\}$. Since $m$ is big and
$\kappa_2>0$,
\[
2\ell_m>2\pi_m+2\kappa=2\pi_m+2\kappa_2+2\log3>2\pi_m+\log(p-m)+\kappa_2 ,
\]
so the first alternative of (F) fails. Hence $\ell_p<\ell_p+t_p<2\pi_m+\kappa_2$.

(b) Let $p\in\{m+1,m+2,m+3\}$. Since $\pi$ is increasing and $m$ is big,
$\pi_p\ge\pi_{m+1}=\pi_m+\ell_m>2\pi_m+\kappa$. If $p$ were big, then
$\ell_p>\pi_p+\kappa>2\pi_m+2\kappa>2\pi_m+\kappa_2$, contradicting (a).
\end{proof}

\begin{lemma}[Evolution of $v_n$]\label{lem:v}
\begin{enumerate}[label=\textup{(\roman*)}]
\item If $m$ is not big, then $v_{m+1}\le v_m+\kappa\,2^{-m-1}$.
\item If $m$ is big, then for $j=1,2,3,4$,
\[
v_{m+j}<\frac{1+2j}{2^j}\,v_m+j\kappa_2\,2^{-m-j};
\]
in particular $v_{m+4}<\tfrac{9}{16}v_m+\kappa_2\,2^{-m-2}$ and
$v_{m+j}<\tfrac32v_m+\kappa_2\,2^{-m}$ for $j=1,2,3$.
\end{enumerate}
\end{lemma}

\begin{proof}
(i) $\pi_{m+1}=\pi_m+\ell_m\le2\pi_m+\kappa$. (ii) By Lemma~\ref{lem:frozen}(a),
$\pi_{m+j}=\pi_m+\sum_{i=0}^{j-1}\ell_{m+i}<(1+2j)\pi_m+j\kappa_2$ for $1\le j\le4$.
Divide by $2^{m+j}$; the coefficients $(1+2j)/2^j$ are $3/2,\,5/4,\,7/8,\,9/16$, and
$v_m\ge0$.
\end{proof}

\subsection{Case I: infinitely many big indices}
Let $B$ be the set of big indices and
$I:=\bigcup_{m\in B}\{m+1,m+2,m+3\}$. By Lemma~\ref{lem:frozen}(b), $I\cap B=\emptyset$.
Enumerate the complement, the set of \emph{checkpoints},
$\N_{\ge1}\setminus I=\{c_1<c_2<\cdots\}$.

\begin{lemma}\label{lem:checkpoints}
If $c_j\in B$ then $c_{j+1}=c_j+4$; if $c_j\notin B$ then $c_{j+1}=c_j+1$.
\end{lemma}

\begin{proof}
If $c_j\in B$, then $c_j+1,c_j+2,c_j+3\in I$. Moreover $c_j+4\notin I$, since otherwise
one of $c_j+1,c_j+2,c_j+3$ would be big, contradicting Lemma~\ref{lem:frozen}(b).
If $c_j\notin B$ and $c_j+1\in I$, then $c_j+1=m+i$ with $m\in B$, $i\in\{1,2,3\}$;
since $c_j\notin B$ we have $m\ne c_j$, so $m\in\{c_j-2,c_j-1\}$ and then $c_j\in I$,
a contradiction.
\end{proof}

Put $w_j:=v_{c_j}$ and $\varepsilon_j:=\kappa\,2^{-c_j}$, so that
$\sum_j\varepsilon_j<\infty$. By Lemmas~\ref{lem:v} and~\ref{lem:checkpoints},
\begin{equation}\label{eq:w}
w_{j+1}<\tfrac9{16}w_j+\varepsilon_j\ \ (c_j\in B),\qquad
w_{j+1}<w_j+\varepsilon_j\ \ (c_j\notin B).
\end{equation}
In particular $w_j\le w_1+\sum_i\varepsilon_i$, so $W:=\limsup w_j$ is finite. Let
$\eta>0$ and choose $J$ with $w_i<W+\eta$ for $i\ge J$ and $\sum_{i\ge J}\varepsilon_i<\eta$.
Since $c_j\in B$ for infinitely many $j$, every sufficiently large $j'$ has a largest
index $j\in[J,j')$ with $c_j\in B$; iterating \eqref{eq:w} from $j$ to $j'$ gives
\[
w_{j'}<\tfrac9{16}w_j+\sum_{i\ge J}\varepsilon_i<\tfrac9{16}(W+\eta)+\eta .
\]
Hence $W\le\tfrac9{16}(W+\eta)+\eta$ for every $\eta>0$, so $W=0$ and $w_j\to0$.
Finally, every $n\in I$ has the form $n=m+i$ with $m=c_j\in B$ and $i\in\{1,2,3\}$, and
Lemma~\ref{lem:v}(ii) gives $v_n<\tfrac32w_j+\kappa_2 2^{-c_j}\to0$. Thus
$v_n\to0$.

\subsection{Case II: finitely many big indices}
Let $N_0$ exceed every big index. For $n\ge N_0$, Lemma~\ref{lem:v}(i) shows that
$u_n:=v_n+\kappa2^{-n}$ is non-increasing, so $v_n\to V$ for some $V\ge0$. Suppose
$V>0$. Then
\[
\ell_n=\pi_{n+1}-\pi_n=2^n\bigl(2v_{n+1}-v_n\bigr)=2^n\bigl(V+o(1)\bigr),
\qquad \ell_{n+1}-\ell_n=2^n\bigl(V+o(1)\bigr).
\]
Choose $N_1\ge N_0$ with $\ell_{k+1}-\ell_k\ge2^{k-1}V$ for all $k\ge N_1$. For $n\ge N_1$
and $k>n$ we get $\ell_k-\ell_n\ge(k-n)2^{n-1}V$, hence
\[
R_n=e^{-\ell_n}\sum_{k\ge n}e^{-(\ell_k-\ell_n)}\le e^{-\ell_n}\Bigl(1+\sum_{i\ge1}e^{-i2^{n-1}V}\Bigr),
\qquad t_n\ge\ell_n-o(1).
\]
Inserting this into (K) yields
$\ell_n+\ell_{n+1}<2\pi_n+\kappa_1+o(1)$, that is,
\[
2^n(V+o(1))+2^{n+1}(V+o(1))<2\cdot2^n(V+o(1))+\kappa_1+o(1).
\]
Dividing by $2^n$ and letting $n\to\infty$ gives $3V\le2V$, contradicting $V>0$.
Hence $V=0$, i.e.\ $v_n\to0$.

\subsection{Conclusion}
In both cases $\pi_n=o(2^n)$. By (J), $\ell_n<2\pi_n+\kappa_2=o(2^n)$. Finally
$a_n=\alpha_n-h\le(1+|h|)\alpha_n$, so $\log a_n\le\ell_n+\log(1+|h|)=o(2^n)$.
Since $a_n\ge1$ for all large $n$ (the sequence is strictly increasing), this is
equivalent to $a_n^{1/2^n}\to1$. This proves Theorem~\ref{thm:A} for integer shifts
$h<h'$; the case $h>h'$ follows by symmetry.

\subsection{Rational shifts}\label{subsec:rational}
Let $t<t'$ be rational with $\sum1/(a_n+t)$, $\sum1/(a_n+t')\in\Q$. Write
$t=u/w$, $t'=u'/w$ with $w\ge1$ and put $b_n:=wa_n+u$. Then $(b_n)$ is a strictly
increasing integer sequence, $\sum1/b_n=w^{-1}\sum1/(a_n+t)\in\Q$ and
$\sum1/(b_n+(u'-u))=w^{-1}\sum1/(a_n+t')\in\Q$. After discarding finitely many terms,
$b_n\ge1$, and the integer case gives $\log b_n=o(2^n)$. Since $a_n=(b_n-u)/w$,
$\log a_n=o(2^n)$ as well. \qed

\begin{remark}[Affine images]\label{rem:affine}
The argument only uses that the two denominators are affine functions of $a_n$. Let
$f(x)=Ux+V$ and $g(x)=U'x+V'$ with integers $U,U'\ge1$ and $\Delta^*:=U'V-UV'\neq0$
(i.e.\ $f,g$ are not proportional), and suppose that $\sum1/f(a_n)$ and
$\sum1/g(a_n)$ are rational. With $\alpha_k=f(a_k)$, $\beta_k=g(a_k)$ one has
$U'\alpha_k-U\beta_k=\Delta^*$, hence
$1/(\alpha_k\beta_k)=(U'/\beta_k-U/\alpha_k)/\Delta^*$, and Lemmas~\ref{lem:identities}--%
\ref{lem:cross} hold verbatim with $D_n=(U'R'_n-UR_n)/\Delta^*$ and $|\Delta^*|$ in place
of $\Delta$ (for $U=U'=1$ this is the translate case, with $\Delta^*=-\Delta$). In
Lemma~\ref{lem:size} the factors $2$ and $e^{\Delta}$ are replaced by constants
depending on $U,U',V,V'$, except that now $\log\Pi'_n=\log\Pi_n+n\log(U'/U)+O(1)$.
Consequently (K), (F), (J) hold with $\kappa_1,\kappa_2$ replaced by $\kappa_i+O(n)$
(resp.\ $O(m)$); defining $m$ to be big if $\ell_m>\pi_m+\kappa+Cm$ for a suitable $C$,
the proof of Sections~3.1--3.4 goes through unchanged, because all error terms remain
summable after division by $2^n$ and are $o(2^n)$ in Case~II. Hence
\emph{$a_n^{1/2^n}\to1$ whenever two non-proportional rational affine images of $(a_n)$
have rational reciprocal sums.}
\end{remark}

\begin{remark}[What the proof uses]\label{rem:uses}
Case~II uses the cross-tail inequality (K) together with the definition of big
indices; Case~I uses only the frozen-jump dichotomy (F) and its special case (J), which
is also used in the conclusion. The single-tail integrality (Lemma~\ref{lem:single})
enters only through Lemmas~\ref{lem:defect} and~\ref{lem:cross}. Neither of the two
rational series alone would suffice: for one series Sylvester's sequence satisfies
$\lim a_n^{1/2^n}>1$. Strict monotonicity can be weakened to monotonicity of $(\alpha_k)$.
\end{remark}

%======================================================================
\section{Multilinear cross-tail integrality and Theorem~\ref{thm:B}}\label{sec:B}
%======================================================================

Let $d\ge1$, $h_1<\dots<h_d$ be integers and let $(a_n)$ be as in Theorem~\ref{thm:B}.
Write $\alpha^{(i)}_k:=a_k+h_i$, $S_i:=\sum_k1/\alpha^{(i)}_k\in\Q$, fix $Q\ge1$ with
$QS_i\in\Z$ for all $i$, and put
\[
\Pi^{(i)}_n:=\prod_{k<n}\alpha^{(i)}_k,\qquad R^{(i)}_n:=\sum_{k\ge n}\frac1{\alpha^{(i)}_k},
\qquad
O_n:=\sum_{\substack{k_1,\dots,k_d\ge n\\ \text{pairwise distinct}}}\ \prod_{i=1}^d\frac1{\alpha^{(i)}_{k_i}} .
\]
For $d=1$, $O_n=R^{(1)}_n$; for $d=2$, $O_n=E_n$ (with $\alpha=\alpha^{(1)}$,
$\beta=\alpha^{(2)}$).

For a finite set $B\subseteq\{1,\dots,d\}$ and $i\in B$ put
$c_{B,i}:=\prod_{j\in B\setminus\{i\}}(h_j-h_i)^{-1}$, so that the partial fraction
decomposition
\begin{equation}\label{eq:pf}
\prod_{i\in B}\frac1{x+h_i}=\sum_{i\in B}\frac{c_{B,i}}{x+h_i}
\end{equation}
holds. Let $\mathcal P_d$ be the set of partitions of $\{1,\dots,d\}$ and, for
$\sigma\in\mathcal P_d$, $\mu(\sigma):=\prod_{B\in\sigma}(-1)^{|B|-1}(|B|-1)!$, the
M\"obius function $\mu(\hat0,\sigma)$ of the partition lattice.

\begin{lemma}[M\"obius expansion]\label{lem:moebius}
For every $n\ge1$,
\[
O_n=\sum_{\sigma\in\mathcal P_d}\mu(\sigma)\prod_{B\in\sigma}\ \sum_{i\in B}c_{B,i}\,R^{(i)}_n .
\]
\end{lemma}

\begin{proof}
For a tuple $\mathbf k=(k_1,\dots,k_d)$ let $\ker\mathbf k\in\mathcal P_d$ be the partition
of $\{1,\dots,d\}$ into the classes of equal coordinates. For $\sigma\in\mathcal P_d$,
\[
F(\sigma):=\prod_{B\in\sigma}\ \sum_{k\ge n}\prod_{i\in B}\frac1{\alpha_k^{(i)}}
=\sum_{\mathbf k:\ \ker\mathbf k\ \ge\ \sigma}\ \prod_{i=1}^d\frac1{\alpha^{(i)}_{k_i}},
\]
where $\ge$ is the refinement order (all series converge absolutely). By M\"obius
inversion on $\mathcal P_d$, the sum over tuples with $\ker\mathbf k=\hat0$ (pairwise
distinct coordinates) equals $\sum_\sigma\mu(\hat0,\sigma)F(\sigma)$. Finally
$\sum_{k\ge n}\prod_{i\in B}1/\alpha^{(i)}_k=\sum_{i\in B}c_{B,i}R^{(i)}_n$ by
\eqref{eq:pf}.
\end{proof}

\begin{lemma}[Multilinear cross-tail integrality]\label{lem:multi}
Let $\mathcal D:=\prod_{1\le i<j\le d}(h_j-h_i)$. For every $n\ge1$,
\[
T_n:=\mathcal D\,Q^d\prod_{i=1}^d\Pi^{(i)}_n\cdot O_n
\]
is a positive integer.
\end{lemma}

\begin{proof}
Expanding the products in Lemma~\ref{lem:moebius}, $O_n$ is a sum of terms
$\mu(\sigma)\prod_{B\in\sigma}c_{B,i_B}R^{(i_B)}_n$ with $i_B\in B$; the
representatives $i_B$ are distinct because the blocks are disjoint. Multiplying such a
term by $Q^d\prod_i\Pi^{(i)}_n$ gives
\[
\mu(\sigma)\prod_{B\in\sigma}c_{B,i_B}\cdot\prod_{B\in\sigma}\bigl(Q\Pi^{(i_B)}_nR^{(i_B)}_n\bigr)
\cdot\prod_{i\notin\{i_B\}}Q\Pi^{(i)}_n .
\]
The middle factors are integers by the argument of Lemma~\ref{lem:single}. The
denominator of $\prod_Bc_{B,i_B}$ is
$\prod_B\prod_{j\in B\setminus\{i_B\}}|h_j-h_{i_B}|$, a product over pairwise
distinct unordered pairs $\{i_B,j\}$, which divides $|\mathcal D|$. Hence
$T_n\in\Z$, and $T_n>0$ since $O_n>0$.
\end{proof}

\begin{lemma}[Size bound]\label{lem:multisize}
Put $G_m:=\sum_{k\ge m}1/\alpha^{(1)}_k$. Then $O_n\le d!\prod_{j=0}^{d-1}G_{n+j}$.
\end{lemma}

\begin{proof}
Since $1/\alpha^{(i)}_k\le1/\alpha^{(1)}_k$, symmetrization gives
$O_n\le d!\sum_{n\le k_1<\dots<k_d}\prod_i 1/\alpha^{(1)}_{k_i}$, and relaxing
$k_1<\dots<k_d$ to the independent conditions $k_j\ge n+j-1$ only enlarges the sum.
\end{proof}

\begin{proof}[Proof of Theorem~\ref{thm:B}]
We first note that \eqref{eq:regular} implies \eqref{eq:tails}. Indeed, if
$a_{k+1}\ge\lambda a_k$ for all $k\ge k_0$, where $\lambda>1$, then
\[
a_n\sum_{k\ge n}\frac1{a_k}\le\frac{\lambda}{\lambda-1}\quad(n\ge k_0),\qquad
a_n\sum_{k\ge n}\frac1{a_k}\le a_{k_0}\sum_{k\ge1}\frac1{a_k}\quad(n<k_0).
\]

Write $\alpha_k:=\alpha^{(1)}_k$, $L_k:=\log\alpha_k$ and $P_n:=\sum_{k<n}L_k$. Since
$\alpha_k=a_k+h_1\ge1$, we have $\alpha_k\le(1+|h_1|)a_k$, and $\alpha_k\ge a_k/2$
whenever $a_k\ge2|h_1|$. Hence \eqref{eq:tails} yields a constant $C_0\ge1$ with
$G_m=\sum_{k\ge m}1/\alpha_k\le C_0/\alpha_m$ for all $m\ge1$ (for the finitely many
$k$ with $a_k<2|h_1|$ one uses $\alpha_k\ge1$ and enlarges $C_0$).
Moreover $\Pi^{(i)}_n\le\Pi^{(1)}_n\prod_k(1+(h_d-h_1)/\alpha_k)\le C_1\Pi^{(1)}_n$ with
$C_1:=\exp((h_d-h_1)S_1)$. Lemmas~\ref{lem:multi} and~\ref{lem:multisize} give
\[
1\le T_n\le |\mathcal D|\,Q^dC_1^d\,(\Pi^{(1)}_n)^d\cdot d!\,C_0^d\prod_{j=0}^{d-1}\alpha_{n+j}^{-1},
\]
that is, $\sum_{j<d}L_{n+j}\le dP_n+C$ with
$C:=\log(|\mathcal D|Q^dC_1^dd!C_0^d)\ge0$. Equivalently $P_{n+d}\le(d+1)P_n+C$, so
$\tilde P_n:=P_n+C/d$ satisfies $\tilde P_{n+d}\le(d+1)\tilde P_n$ and therefore
$\tilde P_n\le C'(d+1)^{n/d}$. Finally $L_n\le P_{n+1}$ and
$\log a_n\le L_n+\log(1+|h_1|)$. Since the $a_k$ are positive integers and
$\sum_k|\log a_k-L_k|<\infty$ (because $|\log(1+h_1/a_k)|=O(1/a_k)$ for large $k$ and
$\sum1/a_k<\infty$), the displayed inequality also holds with $\log a_k$ in place of
$L_k$ after enlarging $C$.
\end{proof}

\begin{remark}\label{rem:Bsharp}
For $d=1$ the recursion $P_{n+1}\le2P_n+C$ is the folklore bound and is attained by
Sylvester-type sequences, so Theorem~\ref{thm:B} is sharp for $d=1$. For $d=2$,
Theorem~\ref{thm:C} gives regular sequences with every exponent $\beta<\beta_*$, so the
optimal regular exponent lies in $[\beta_*,\sqrt3]$. For general $d$, \cite{KT25} gives
$\beta<((2d+2)/(2d+1))^{1/d}$ (their construction places one term in each of a sequence
of rapidly growing windows, so it satisfies \eqref{eq:regular}), and
Theorem~\ref{thm:B} gives $\beta\le(d+1)^{1/d}$.
\end{remark}

%======================================================================
\section{The construction: proof of Theorem~\ref{thm:C}}\label{sec:C}
%======================================================================

\subsection{Two targets}
We construct the sequence term by term, prescribing the two sums in advance. For real
numbers $\tau,\delta$ put
\[
\Gamma:=\tau^2-\delta,\qquad s:=\frac{\Gamma}{2\tau}\quad(\tau>0),
\]
and for an integer $a\ge2$ with $x:=1/a$ put
\[
\tau^{a}:=\tau-\frac1a,\qquad \delta^{a}:=\delta-\frac1{a(a-1)},\qquad
\Gamma^{a}:=(\tau^{a})^2-\delta^{a},\qquad s^{a}:=\frac{\Gamma^{a}}{2\tau^{a}} .
\]
Since $1/(a(a-1))=x^2/(1-x)$, a direct computation gives the \emph{one-step identity}
\begin{equation}\label{eq:onestep}
\Gamma^{a}=\Gamma-2x\,\tau^{a}+\frac{x^3}{1-x}.
\end{equation}
(Indeed, $\Gamma^a-\Gamma=-2\tau x+x^2+x^2/(1-x)$ and $-2\tau x=-2x\tau^a-2x^2$.)
Heuristically, if $\tau$ and $\delta$ are the tails $\sum_{k\ge n}1/a_k$ and
$\sum_{k\ge n}1/(a_k(a_k-1))$ of a rapidly growing sequence, then
$\Gamma=2\sum_{n\le j<k}\frac1{a_ja_k}-\sum_{k\ge n}\frac{1}{a_k^2(a_k-1)}\approx
\frac{2}{a_na_{n+1}}$, so that $s\approx 1/a_{n+1}$ is the size of the \emph{next}
tail, while $\tau\approx1/a_n$.

\begin{lemma}[One step]\label{lem:step}
Let $K\ge64$. Let $\tau,\delta$ be real numbers with $0<\tau\le K^{-2}$ and suppose that
$\Gamma=\tau^2-\delta$ and $s=\Gamma/(2\tau)$ satisfy
\begin{equation}\label{eq:inv}
K\tau^2\le s\le \tau/K .
\end{equation}
Let $y$ be a real number with
\begin{equation}\label{eq:ychoice}
\frac{8\tau^3}{s}\le y\le\frac{s}{4K}.
\end{equation}
Then there is an integer $a>1/\tau$ such that:
\begin{enumerate}[label=\textup{(\alph*)}]
\item $\tau^a>0$, $\delta^a>0$ and $\Gamma^a>0$;
\item $(1-4/K)\,s\le\tau^a\le(1+6/K)\,s$;
\item $y\le s^a\le 2y$;
\item $\tau^a\le K^{-2}$ and $K(\tau^a)^2\le s^a\le\tau^a/K$;
\item $1/\tau<a\le(1+4/K)/\tau$.
\end{enumerate}
\end{lemma}

\begin{proof}
For an integer $a>1/\tau$ the number $t:=\tau^a=\tau-1/a$ lies in $(0,\tau)$ and
$x=\tau-t$. By \eqref{eq:onestep}, $\Gamma^a=\Gamma-g(t)$, where
\[
g(t):=2(\tau-t)t-\frac{(\tau-t)^3}{1-\tau+t}\qquad(0\le t<\tau).
\]

\emph{Step 1: slope of $g$.} We have $g'(t)=2\tau-4t+h(t)$ with
\[
h(t)=\frac{3(\tau-t)^2(1-\tau+t)+(\tau-t)^3}{(1-\tau+t)^2}\in\Bigl[0,\frac{4\tau^2}{(1-\tau)^2}\Bigr]
\subseteq[0,5\tau^2].
\]
Hence, for $0\le t\le 2\tau/K$, using $\tau\le K^{-2}\le K^{-1}$,
\begin{equation}\label{eq:slope}
2\tau(1-4/K)\le g'(t)\le 2\tau+5\tau^2\le 2\tau(1+3/K).
\end{equation}
Also $g(0)=-\tau^3/(1-\tau)\in[-2\tau^3,0)$.

\emph{Step 2: the root $t_0$.} By \eqref{eq:slope},
$g(2\tau/K)\ge g(0)+\frac{4\tau^2}{K}(1-\frac4K)\ge\frac{2\tau^2}{K}\bigl(2-\frac8K-K\tau\bigr)
>\frac{2\tau^2}{K}\ge2\tau s=\Gamma$, since $K\tau\le 1/K$ and $K\ge64$. As $g$ is
increasing on $[0,2\tau/K]$ and $g(0)<0<\Gamma$, there is a unique $t_0\in(0,2\tau/K)$
with $g(t_0)=\Gamma$. By the mean value theorem and \eqref{eq:slope},
$2\tau(1-4/K)t_0\le\Gamma-g(0)\le 2\tau(1+3/K)t_0$, and
$2\tau s=\Gamma\le\Gamma-g(0)\le 2\tau(s+\tau^2)$. Using $\tau^2\le s/K$ from
\eqref{eq:inv} we get
\begin{equation}\label{eq:t0}
(1-3/K)\,s\le\frac{s}{1+3/K}\le t_0\le\frac{s+\tau^2}{1-4/K}\le\frac{1+1/K}{1-4/K}\,s\le(1+6/K)\,s .
\end{equation}

\emph{Step 3: the target interval.} For $c>0$ the function $t\mapsto\Gamma-g(t)-2ct$ is
strictly decreasing on $[0,t_0]$, positive at $t=0$ and negative at $t=t_0$; let
$t_c\in(0,t_0)$ be its zero. For $t\in(0,t_0]$ we then have
$(\Gamma-g(t))/(2t)\ge c$ if and only if $t\le t_c$. Consequently, if $t=\tau^a\in(0,t_0]$,
\begin{equation}\label{eq:window}
y\le s^a\le2y\quad\Longleftrightarrow\quad t\in[t_{2y},t_y].
\end{equation}
We estimate this interval. Since $\Gamma-g(t_{2y})=4yt_{2y}$ and $\Gamma-g(t_0)=0$,
\eqref{eq:slope} gives $t_0-t_{2y}\le 4yt_{2y}/(2\tau(1-4/K))\le 2yt_0/(\tau(1-4/K))$.
By \eqref{eq:ychoice} and \eqref{eq:inv}, $y\le s/(4K)\le\tau/(4K^2)$, hence
$t_0-t_{2y}\le t_0/K^2$ and, by \eqref{eq:t0},
\begin{equation}\label{eq:t2y}
t_{2y}\ge(1-K^{-2})(1-3/K)\,s\ge(1-4/K)\,s .
\end{equation}
Next, the function $F(t):=\Gamma-g(t)-2yt$ satisfies $F(t_{2y})=4yt_{2y}-2yt_{2y}=2yt_{2y}$,
$F(t_y)=0$ and $|F'|\le 2\tau(1+3/K)+2y\le2\tau(1+4/K)$ on $[0,t_0]$. Therefore, using
\eqref{eq:t2y}, $K\ge64$ and the lower bound in \eqref{eq:ychoice},
\[
t_y-t_{2y}\ge\frac{2yt_{2y}}{2\tau(1+4/K)}\ge\frac{1-4/K}{1+4/K}\cdot\frac{ys}{\tau}
\ge\frac{ys}{2\tau}\ge4\tau^2 .
\]

\emph{Step 4: choice of $a$.} The numbers $\tau-1/a$ with integers $a>1/\tau$ increase
to $\tau$; consecutive ones differ by $1/(a(a+1))<1/a^2<\tau^2$, and the smallest one
is $\le\tau-\tau/(1+\tau)<\tau^2$. The interval $[t_{2y},t_y]$ has length $\ge4\tau^2$ and
lies in $[\tau^2,\tau)$, because $t_{2y}\ge(1-4/K)K\tau^2\ge\tau^2$ by \eqref{eq:t2y},
\eqref{eq:inv}, and $t_y<t_0<2\tau/K<\tau$. Hence it contains $\tau-1/a$ for some
integer $a>1/\tau$; we fix such an $a$.

\emph{Step 5: verification.} (c) holds by \eqref{eq:window}. (b) follows from
$t_{2y}\le\tau^a\le t_y<t_0$, \eqref{eq:t2y} and \eqref{eq:t0}. For (a): $\tau^a>0$,
$\Gamma^a=2\tau^as^a>0$, and $\delta^a=(\tau^a)^2-\Gamma^a=\tau^a(\tau^a-2s^a)>0$ because
$2s^a\le4y\le s/K<(1-4/K)s\le\tau^a$. For (d): $\tau^a\le(1+6/K)s\le(1+6/K)\tau/K\le\tau\le
K^{-2}$; $s^a\le2y\le s/(2K)\le\tau^a/(2K(1-4/K))\le\tau^a/K$; and
$K(\tau^a)^2\le K(1+6/K)^2s^2\le 2Ks^2\le 8\tau^3/s\le y\le s^a$, where
$2Ks^2\le8\tau^3/s$ holds because $s\le\tau/K$ gives $Ks^3\le\tau^3/K^2\le4\tau^3$.
For (e): $a=1/(\tau-\tau^a)$ and
$\tau-\tau^a\ge\tau-(1+6/K)\tau/K\ge(1-2/K)\tau$, so $a\le\tau^{-1}/(1-2/K)\le(1+4/K)/\tau$;
and $a>1/\tau$ since $\tau^a>0$.
\end{proof}

\subsection{Proof of Theorem~\ref{thm:C}}
Fix $\beta\in(1,\beta_*)$, so that $\beta^2+\beta<3$, and put $K:=64$. Let $N$ be a
large integer (conditions below), $\tau_1:=1/N$, choose a rational number $s_1$ with
$\tfrac12N^{-\beta}\le s_1\le N^{-\beta}$, and put $\delta_1:=\tau_1^2-2\tau_1s_1\in\Q$.
Recursively, given $\tau_n,\delta_n$, let $\Gamma_n:=\tau_n^2-\delta_n$,
$s_n:=\Gamma_n/(2\tau_n)$, $y_n:=s_n^{\beta}$, let $a_n$ be the integer provided by
Lemma~\ref{lem:step} for $(\tau,\delta,y)=(\tau_n,\delta_n,y_n)$, and put
\[
\tau_{n+1}:=\tau_n-\frac1{a_n},\qquad \delta_{n+1}:=\delta_n-\frac1{a_n(a_n-1)} .
\]
Write $\lambda_n:=-\log\tau_n$ and $\nu_n:=-\log s_n$. We show by induction that for all
$n\ge1$ the lemma is applicable and
\[
(\mathrm H_n)\qquad \tau_n\le K^{-2},\qquad K\tau_n^2\le s_n\le\tau_n/K,\qquad
|\nu_n-\beta\lambda_n|\le1,\qquad \lambda_n\ge\lambda_1 .
\]
\emph{Conditions on $N$.} With $\lambda_1=\log N$ we require
\begin{equation}\label{eq:Nlarge}
\lambda_1\ge\max\Bigl\{2\log K,\ \frac{\log K}{\beta-1},\ \frac{\log(2K)}{2-\beta},\
\frac{1+\log(4K)/(\beta-1)}{\beta},\ \frac{\beta+1+\log8}{3-\beta-\beta^2},\
\frac{1.1}{\beta-1}\Bigr\}.
\end{equation}

\emph{Base case.} $\tau_1=1/N\le K^{-2}$; $\nu_1\in[\beta\lambda_1,\beta\lambda_1+\log2]$,
so $|\nu_1-\beta\lambda_1|\le1$; and $K\tau_1^2\le s_1\le\tau_1/K$ is equivalent to
$\lambda_1+\log K\le\nu_1\le2\lambda_1-\log K$, which follows from \eqref{eq:Nlarge}.

\emph{Applicability.} Assume $(\mathrm H_n)$. Condition \eqref{eq:ychoice} for
$y_n=s_n^\beta$ reads $(\beta-1)\nu_n\ge\log(4K)$ and $(\beta+1)\nu_n\le3\lambda_n-\log8$.
The first follows from $\nu_n\ge\beta\lambda_n-1\ge\beta\lambda_1-1$ and
\eqref{eq:Nlarge}. For the second,
$(\beta+1)\nu_n\le(\beta+1)(\beta\lambda_n+1)=(\beta^2+\beta)\lambda_n+\beta+1\le3\lambda_n-\log8$,
because $(3-\beta-\beta^2)\lambda_n\ge(3-\beta-\beta^2)\lambda_1\ge\beta+1+\log8$.

\emph{Induction step.} Lemma~\ref{lem:step}(d) gives the first two parts of
$(\mathrm H_{n+1})$. By (b) and (c), with $\log(1+6/64)<0.09$ and $-\log(1-4/64)<0.07$,
\[
\nu_n-0.09\le\lambda_{n+1}\le\nu_n+0.07,\qquad \beta\nu_n-\log2\le\nu_{n+1}\le\beta\nu_n .
\]
Hence $\nu_{n+1}-\beta\lambda_{n+1}\in[-\log2-0.07\beta,\,0.09\beta]\subseteq[-1,1]$, and
$\lambda_{n+1}\ge\nu_n-0.09\ge\beta\lambda_n-1.09\ge\lambda_n$ by \eqref{eq:Nlarge}.
This proves $(\mathrm H_{n+1})$.

\emph{Monotonicity.} By (e), $1/\tau_n<a_n\le(1+4/K)/\tau_n$, and by (b) and
$(\mathrm H_n)$, $\tau_{n+1}\le(1+6/K)s_n\le(1+6/K)\tau_n/K$. Hence
$a_{n+1}>1/\tau_{n+1}\ge K/((1+6/K)\tau_n)>(1+4/K)/\tau_n\ge a_n$, and $a_1>N\ge2$.
Moreover $a_{n+1}/a_n\ge e^{\lambda_{n+1}-\lambda_n}/(1+4/K)$ and
$\lambda_{n+1}-\lambda_n\ge(\beta-1)\lambda_n-1.09$. With $c:=1.09/(\beta-1)$ the
inequality $\lambda_{n+1}\ge\beta\lambda_n-1.09$ gives
$\lambda_n-c\ge\beta^{n-1}(\lambda_1-c)$, and $\lambda_1>c$ by \eqref{eq:Nlarge}; hence
$\lambda_n\to\infty$ and $a_{n+1}/a_n\to\infty$.

\emph{Rationality.} We have $\tau_{n+1}\le\tau_n/2$, so $\tau_n\to0$, and
$0<\delta_n=\tau_n^2-\Gamma_n<\tau_n^2$ (for $n\ge2$ by Lemma~\ref{lem:step}(a), and for
$n=1$ because $\delta_1=\tau_1(\tau_1-2s_1)$ with $2s_1<\tau_1$), so $\delta_n\to0$.
Since $\tau_n=\tau_1-\sum_{k<n}1/a_k$ and $\delta_n=\delta_1-\sum_{k<n}1/(a_k(a_k-1))$,
\[
\sum_{k\ge1}\frac1{a_k}=\tau_1\in\Q,\qquad
\sum_{k\ge1}\frac1{a_k-1}=\sum_{k\ge1}\Bigl(\frac1{a_k}+\frac1{a_k(a_k-1)}\Bigr)=\tau_1+\delta_1\in\Q .
\]

\emph{Growth.} Since $\nu_{n+1}-\beta\nu_n\in[-\log2,0]$, the limit
\[
C:=\lim_{n\to\infty}\frac{\nu_n}{\beta^n}=\frac{\nu_1}{\beta}+\sum_{k\ge1}\frac{\nu_{k+1}-\beta\nu_k}{\beta^{k+1}}
\ \ge\ \frac{\nu_1}{\beta}-\frac{\log2}{\beta(\beta-1)}
\]
exists and is positive by \eqref{eq:Nlarge}. As
$|\log a_{n+1}-\lambda_{n+1}|\le\log(1+4/K)<0.07$ and $|\lambda_{n+1}-\nu_n|\le0.09$, we
get $\log a_{n+1}/\beta^{n+1}\to C/\beta$. This proves Theorem~\ref{thm:C} with
$c=C/\beta$. \qed

\begin{remark}[Where $\beta_*$ comes from]\label{rem:heuristic}
The only constraint in the construction is \eqref{eq:ychoice}: as $a$ runs through the
integers, $\Gamma^a$ moves in steps of size $\asymp\tau^3\asymp a_n^{-3}$
(by \eqref{eq:slope}), while $\Gamma_{n+1}=2\tau_{n+1}s_{n+1}\asymp
1/(a_{n+1}a_{n+2})$. So one needs $a_{n+1}a_{n+2}\lesssim a_n^3$, i.e.\
$\beta+\beta^2\le3$ for $\log a_n\asymp\beta^n$, while the window \eqref{eq:ychoice}
offers $\asymp a_n^3/(a_{n+1}a_{n+2})$ admissible choices at step $n$. Treating the
fractional positions of the targets as independent and uniformly distributed, the
expected number of sequences that follow a growth profile $\log a_k=L_k$ up to step $n$ is
therefore
$\exp\bigl(\sum_{k\le n}(3L_k-L_{k+1}-L_{k+2})+O(n)\bigr)$, which tends to $0$ doubly
exponentially if $L_k\sim c\beta^k$ with $\beta>\beta_*$ (and along a subsequence if
only $L_k\asymp\beta^k$). This heuristic only models sequences in which each term
nearly exhausts its tail. Sequences satisfying \eqref{eq:regular} may also contain blocks
of terms of comparable size, so the heuristic merely suggests that
$\beta_2^{\mathrm{reg}}=\beta_*$. By contrast, the integrality bound of
Theorem~\ref{thm:B} corresponds to the extreme situation in which the cross-tail
integers $M_n$ of Lemma~\ref{lem:cross} remain bounded, which is what growth with
exponent $\sqrt3$ would require.
\end{remark}

\begin{remark}[Other distances]\label{rem:otherDelta}
For two integer translates at distance $\Delta\ge1$ one works with $x=1/\beta_k$,
$\tau=\sum1/\beta_k$ and $\delta=\sum1/(\alpha_k\beta_k)=\sum x^2/(1-\Delta x)$, and
\eqref{eq:onestep} becomes $\Gamma^a=\Gamma-2x\tau^a+\Delta x^3/(1-\Delta x)$. The proof
of Lemma~\ref{lem:step} goes through with $K$ depending on $\Delta$, and gives the same
exponent $\beta_*$; we omit the details.
\end{remark}

%======================================================================
\section{Consequences}\label{sec:cor}
%======================================================================

\begin{proof}[Proof of Corollary~\ref{cor:D}]
The first assertion is Theorem~\ref{thm:A} with $t=-1$, $t'=0$ (and is Kitamura's
theorem \cite{Kit26}); the second is Theorem~\ref{thm:B} with $d=2$ and
$(h_1,h_2)=(-1,0)$. For the third, given $\beta<\beta_*$ pick
$\beta'\in(\max(\beta,1),\beta_*)$ and apply Theorem~\ref{thm:C} with $\beta'$: then
$\log a_n/\beta^n=(\log a_n/\beta'^n)(\beta'/\beta)^n\to\infty$. For the bounds on
$\beta^{\mathrm{reg}}_2$, which is defined with the shifts $\{0,1\}$, substitute
$b_n:=a_n-1$: then $\sum1/b_n$ and $\sum1/(b_n+1)$ are rational, $b_n\ge1$,
$b_{n+1}/b_n\to\infty$ and $\log b_n\sim\log a_n$, so the sequences of
Theorem~\ref{thm:C} give $\beta^{\mathrm{reg}}_2\ge\beta_*$. Conversely, after discarding
finitely many non-positive terms, Theorem~\ref{thm:B} with $(h_1,h_2)=(0,1)$ gives
$\log b_n=O(3^{n/2})$ for every regular sequence, so $b_n^{1/\beta^n}\to\infty$ is
impossible for $\beta\ge\sqrt3$, and $\beta_2^{\mathrm{reg}}\le\sqrt3$.
\end{proof}

\begin{proof}[Proof of Corollary~\ref{cor:E}]
If two distinct rational translates had rational reciprocal sums, Theorem~\ref{thm:A}
would give $a_n^{1/2^n}\to1$, contradicting $\limsup a_n^{1/2^n}>1$.
\end{proof}

\begin{example}[Sylvester's sequence]\label{ex:sylvester}
Sylvester's sequence satisfies $\sum1/s_n=1$ and $s_n^{1/2^n}\to E^2>1$, so by
Corollary~\ref{cor:E}, $\sum_n1/(s_n+t)\notin\Q$ for every rational $t\ne0$ with
$s_n+t\neq0$ for all $n$. This is \emph{not} new: writing $t=u/w$ and $u_n:=ws_n+u$, one
has $u_{n+1}=u_n^2/w+O(u_n)$ and
\[
u_{n+1}-\Bigl(\frac{u_n^2}{w}-u_n+w\Bigr)=-u\Bigl(2s_n+\frac uw-2\Bigr)\neq0
\quad(n\text{ large}),
\]
so the irrationality follows from Duverney's criterion \cite[Cor.~3.4]{Duverney01}; for
integer $t$ it can also be deduced from the Erd\H{o}s--Straus criterion \cite{ErSt64},
and for negative integers $t$ from Badea's criterion \cite{Badea87}. Corollary~\ref{cor:E} requires no recursion.
\end{example}

\begin{corollary}[Infinitely many rational translates]\label{cor:many}
Let $a_1<a_2<\cdots$ be positive integers satisfying \eqref{eq:tails}, and suppose that
$\sum_n1/(a_n+h)\in\Q$ for infinitely many integers $h\ge0$. Then $\log a_n=O(\beta^n)$
for every $\beta>1$.
\end{corollary}

\begin{proof}
Apply Theorem~\ref{thm:B} to $d$ of these shifts and note $(d+1)^{1/d}\to1$.
\end{proof}

\begin{remark}[Related Erd\H{o}s problems]\label{rem:related}
(1) \emph{Problem \#266} asked whether, for every convergent $\sum1/a_n$, some integer
$t\ge1$ makes $\sum1/(a_n+t)$ irrational; Kova\v{c} and Tao \cite{KT25} disproved this
with a strictly increasing sequence for which all such sums are rational. By
Theorem~\ref{thm:A}, every strictly increasing ``universally rational'' sequence satisfies
$a_n^{1/2^n}\to1$, and by Corollary~\ref{cor:many} it satisfies $\log a_n=O(\beta^n)$
for every $\beta>1$ if its tails are dominated by their first terms. For strictly
increasing sequences with $\limsup a_n^{1/2^n}>1$, Corollary~\ref{cor:E} gives the
conclusion of \#266 in the strong form ``all translates but at most one are
irrational''. (For non-monotone sequences these statements fail: rearranging terms does
not change the sums but can make $\limsup a_n^{1/2^n}$ infinite.)

(2) \emph{Problem \#263(i)} asks whether $2^{2^n}$ is an irrationality sequence, i.e.\
whether $\sum1/b_n\notin\Q$ whenever $b_n/2^{2^n}\to1$. Such $(b_n)$ are eventually
strictly increasing with $b_n^{1/2^n}\to2$, so by Corollary~\ref{cor:E} at most one
rational translate of $(b_n)$ has a rational reciprocal sum.

(3) \emph{Problem \#243} concerns sequences with $a_n/a_{n-1}^2\to1$ and rational
$\sum1/a_n$. For such sequences $\log a_n/2^n$ converges to a positive limit, so by
Corollary~\ref{cor:E} no other rational translate of $(a_n)$ has a rational reciprocal
sum.
\end{remark}

%======================================================================
\section{Limits of the method and open problems}\label{sec:limits}
%======================================================================

\subsection{Why the general bound stops at $2^n$}
Theorem~\ref{thm:A} is qualitative, while Theorem~\ref{thm:B} gives an exponential
saving under \eqref{eq:tails}. The proof of Theorem~\ref{thm:A} shows more precisely
that along the checkpoints the normalized quantity $v_n=\pi_n/2^n$ can increase only
by a summable amount, and that it drops by the factor $9/16$ across every frozen
window. What the argument cannot exclude are long \emph{doubling phases}: runs of
non-big indices with $\ell_k=\pi_k+O(1)$, along which $v_n$ stays essentially constant.
By Case~II such a run cannot last forever, and along it the tails cannot be dominated
by their first terms (otherwise (K) would be violated): the remaining mass $R_k$ must
be reserved for a later, very long run of terms of comparable size, which then lowers
$v_n$ drastically. We stress that this is a \emph{heuristic} description of what the
inequalities (K), (F), (J) permit; we do not claim that doubling phases occur for
actual sequences with two rational translates.

There is a precise Diophantine reason why ruling them out is subtle. Put
$y_k:=\alpha_k\beta_k$ and $N_k:=\Delta Q\Pi_k\Pi'_kD_k$ (Lemma~\ref{lem:defect}).
Since $D_{k+1}=D_k-1/y_k$, one has the exact recursions
\[
N_{k+1}=y_kN_k-\Delta Q\,\Pi_k\Pi'_k,\qquad
y_{k+1}N_{k+1}=y_k^2N_k-y_kN_{k+1}+N_{k+2}\qquad(k\ge1).
\]
The second recursion is exactly the one that governs Erd\H{o}s Problem~\#243 for the
rational series $\sum_k1/y_k=D_1$ (with $N_k$ in the role of the integers
$q\prod_{j<k}y_j\sum_{j\ge k}1/y_j$). Along a doubling phase of length at least three
the defect tail $D_k$ is dominated by its first term, so the positive integers $N_k$
stay bounded. Long doubling phases would therefore require a \#243-type behaviour of
the defect series with bounded $N_k$, subject to the additional constraint that every
$y_k$ has the form $\alpha(\alpha+\Delta)$. We do not know how to exclude this.

\begin{remark}[A Pell-type identity]\label{rem:pell}
Put $z_k:=2\alpha_k+\Delta$, so that $4y_k=z_k^2-\Delta^2$. Eliminating $y_k,y_{k+1}$
from the second recursion above gives the exact identity
\[
u_k^2-4N_kN_{k+1}\,z_{k+1}^2=c_k,
\]
where
\[
u_k:=N_kz_k^2-N_k\Delta^2-2N_{k+1},\qquad c_k:=4N_{k+1}^2-4N_kN_{k+1}\Delta^2-16N_kN_{k+2},
\]
for every $k\ge1$ (a direct computation, checked in Appendix~\ref{app:scripts}). Thus
along any run on which the $N_k$ stay bounded, the pairs $(u_k,z_{k+1})$ solve
generalized Pell equations $X^2-4N_kN_{k+1}Y^2=c_k$ with bounded coefficients. If
$N_kN_{k+1}=s^2$ is a perfect square and $u_k+2sz_{k+1}>|c_k|$ (which holds once $z_k$
is large compared with the $N$'s), then the integer $u_k-2sz_{k+1}$ has absolute value
$<1$, hence vanishes, and $c_k=0$, i.e.\ $N_{k+1}(N_{k+1}-\Delta^2N_k)=4N_kN_{k+2}$. In
particular the $N_k$ cannot be constant along a run with large $z_k$, since a constant
value $N$ would give $c_k=-4N^2(3+\Delta^2)\neq0$. Whether this Pell structure can be
iterated along long runs is, in our view, the crux of Question~\ref{q:irregular} below.
\end{remark}

\subsection{Open problems}

\begin{question}\label{q:reg}
Is $\beta_2^{\mathrm{reg}}=\beta_*=(\sqrt{13}-1)/2$? By Theorem~\ref{thm:C} this would
follow if every sequence with $\sum1/a_n,\sum1/(a_n-1)\in\Q$ and
$\liminf a_{n+1}/a_n>1$ satisfied $\log a_n=O(\beta^n)$ for every $\beta>\beta_*$. We
know $\beta_*\le\beta_2^{\mathrm{reg}}\le\sqrt3$.
\end{question}

\begin{question}\label{q:irregular}
Is there $\beta<2$ such that $\log a_n=O(\beta^n)$ for \emph{every} sequence with two
rational translates (without \eqref{eq:tails})? Can irregular sequences beat the regular
exponent?
\end{question}

\begin{question}\label{q:d}
For $d\ge3$ consecutive rational translates, close the gap between
$((2d+2)/(2d+1))^{1/d}$ \cite{KT25} and $(d+1)^{1/d}$ (Theorem~\ref{thm:B}). Does a
controlled $d$-target greedy algorithm improve the lower bound?
\end{question}

The heuristic of Remark~\ref{rem:heuristic} extends formally to $d$ targets: with $d$
prescribed sums the choice of $a_n$ can only steer the highest-order discrepancy, which
moves in steps of size $\asymp a_n^{-(d+1)}$ and has size $\asymp(a_{n+1}\cdots a_{n+d})^{-1}$.
This suggests the threshold $\beta_d^*$ defined by $\beta+\beta^2+\cdots+\beta^d=d+1$.
For $d=1$ this gives $\beta_1^*=2$, the Sylvester threshold; for $d=2$ it gives
$\beta_2^*=\beta_*$; and $\beta_3^*=1.1509\ldots$, $\beta_4^*=1.0912\ldots$, compared with
$((2d+2)/(2d+1))^{1/d}=1.0455\ldots,\ 1.0266\ldots$ from \cite{KT25} and
$(d+1)^{1/d}=1.5874\ldots,\ 1.4953\ldots$ from Theorem~\ref{thm:B}. In a numerical
experiment with three targets ($\sum1/a_k$, $\sum1/(a_k-1)$, $\sum1/(a_k-2)$), where the
choice of $a_n$ steers the triple off-diagonal sum $O_{n+1}$ of Section~\ref{sec:B}, the
largest attainable value of $\log a_{n+3}$ was close to
$4\log a_n-\log a_{n+1}-\log a_{n+2}$, as predicted. The process ran stably for $14$
steps with growth ratio $\log a_{n+1}/\log a_n\approx1.15$. We have not proved a
$d$-target analogue of Lemma~\ref{lem:step} for $d\ge3$, and we do not claim that
$\beta_d^*$ is the correct value of $\beta_d^{\mathrm{reg}}$.

%======================================================================
\section{Comparison with the literature}\label{sec:lit}
%======================================================================

\begin{itemize}[leftmargin=1.4em]
\item The problem appears in \cite[p.~64]{ErGr80}, \cite[p.~334]{Er85},
\cite[pp.~104--105]{Er88} and as \#265 in \cite{Bloom265}. As of October~28, 2026 the
database lists it as open; the page states that it remains open whether
$\limsup a_n^{1/2^n}>1$ can be achieved, and that the remaining question is the precise
exponent.
\item Kova\v{c}'s preprint \cite{Kov24} (arXiv:2406.17593v1) proved that exponential
growth is possible; the later versions of that arXiv entry are the joint paper
\cite{KT25}, which proves doubly exponential growth for $d$ consecutive translates with
any $\beta<((2d+2)/(2d+1))^{1/d}$, i.e.\ $\beta<\sqrt{6/5}$ for two series, and states
that the question beyond the $2^n$ threshold is not answered. Their methods (block
constructions, interval-filling and interior arguments) address the lower bound; we are
not aware of upper-bound arguments of the type of Lemma~\ref{lem:cross} there.
\item \emph{Kitamura \cite{Kit26}.} A Lean~4 formalization (repository
\texttt{KitaKen1/erdos-265-lean}, commit \texttt{16db9f29}, September 2026) proves that
$\log a_n/2^n\to0$ for every strictly increasing sequence with $a_1\ge2$ and
$\sum1/a_n,\sum1/(a_n-1)\in\Q$, hence that $\limsup a_n^{1/2^n}>1$ is impossible. It was
linked from a public comment on the database page dated September~7, 2026, and was
checked independently (with the \emph{Comparator} tool) by a reviewer of the
formal-conjectures pull request \cite{FC6839} in October~2026. Its proof uses the
defect integrality (as in Lemma~\ref{lem:defect}), an envelope $H_n=P_n/\sqrt{D_n}$
satisfying a square recursion $H_{n+1}\le K H_n^2$, and the integrality and the bound
$\le8T_nT_{n+1}$ of the second residual $T_n/(1-T_n)-V_n$, which by
Remark~\ref{rem:residual} equals our cross tail divided by $1-T_n$. (These three lemmas
already imply the case $d=2$ of Theorem~\ref{thm:B} for regular sequences, although
\cite{Kit26} uses them only to exclude a positive normalized limit.) We found no
human-readable account of it. Thus the endpoint for Problem \#265, and the
cross-tail integrality for $\Delta=1$, are due to Kitamura; our contributions are the
extension to arbitrary rational translates and affine images (Theorem~\ref{thm:A},
Remark~\ref{rem:affine}), the frozen-jump/checkpoint route, the multilinear
integrality and Theorem~\ref{thm:B} for $d\ge3$, and the construction of
Theorem~\ref{thm:C}.
\item An AI-generated working report \cite{EPAD265} on \#265 (accessed July 2026) claims
a construction with regular growth exponent $\sqrt{3/2}=1.2247\ldots$ by a fixed-window
block method and argues that such block methods cannot exceed $\sqrt{5/3}=1.2909\ldots$;
we have not verified these claims. Theorem~\ref{thm:C} uses a different, sequential
method and gives every exponent below $\beta_*=1.3027\ldots$.
\item The folklore irrationality criterion ($a_n^{1/2^n}\to\infty$), Badea's criterion
\cite{Badea87}, the Erd\H{o}s--Straus criterion \cite{ErSt64}, Duverney's criteria
\cite{Duverney01} and Koizumi's results \cite{Koizumi25} concern a \emph{single} Ahmes
series. The Sylvester example of Example~\ref{ex:sylvester} already follows from
\cite{Duverney01}.
\item To our knowledge, Theorem~\ref{thm:C}, the multilinear integrality of
Lemma~\ref{lem:multi} and Theorem~\ref{thm:B} for $d\ge3$ are new; a web search up to
October~28, 2026 (including the database, the formal-conjectures project, arXiv, and the
report \cite{EPAD265}) found no construction with exponent above $\sqrt{3/2}$ and no
mention of $\beta_*$. Theorem~\ref{thm:B} for $d=2$ is not stated in \cite{Kit26} but
follows quickly from lemmas proved there, and Theorem~\ref{thm:A} beyond the pair
$(a_n,a_n-1)$ appears to be a routine extension of Kitamura's argument, although we have
not seen it stated. We could not access the live discussion thread of the database
page.
\end{itemize}

\appendix
%======================================================================
\section{Exact verification script}\label{app:scripts}
%======================================================================

The following Python script (standard library only, exact arithmetic via
\texttt{fractions.Fraction}) checks: the identities of Lemma~\ref{lem:identities}; the
integrality and positivity statements of Lemmas~\ref{lem:defect} and~\ref{lem:cross};
the size bounds of Lemma~\ref{lem:size}(i),(ii); the M\"obius expansion
(Lemma~\ref{lem:moebius}), the integrality of $T_n$ (Lemma~\ref{lem:multi}) and the
bound of Lemma~\ref{lem:multisize} for $d=2,3,4$; for Cantor's example
$a_n=\binom n2$ ($n\ge3$), where the tails are known in closed form, the integrality and
positivity of $N_n$ and $M_n$ and the bound of Lemma~\ref{lem:size}(i); the numerical
facts used in the text (the coefficients $(1+2j)/2^j$, Sylvester's telescoping identity,
$E^2=1.59791\ldots$, the Badea comparison and the Duverney comparison of
Example~\ref{ex:sylvester}); the Pell-type identity of Remark~\ref{rem:pell}, both as a
rational identity and on the defect integers $N_k$ of random finite sequences; the
affine version of Remark~\ref{rem:affine}; the residual identity of
Remark~\ref{rem:residual}; and the construction of Section~\ref{sec:C}: it runs the
recursion of the proof of Theorem~\ref{thm:C} (with $\beta=5/4$, $K=64$, $N=10^{11}$)
for eleven steps, with $y_n$ replaced by an $80$-digit rational approximation of
$s_n^\beta$ (Lemma~\ref{lem:step} holds for every real $y$ satisfying \eqref{eq:ychoice}).
At every step it checks, in exact arithmetic, the hypotheses \eqref{eq:inv},
\eqref{eq:ychoice} and all conclusions (a)--(e) of Lemma~\ref{lem:step}, and that the
partial sums reproduce the prescribed rational targets exactly. The inequality
$|\nu_n-\beta\lambda_n|\le1$ is checked in floating point. For a \emph{finite}
increasing sequence the reciprocal sums are automatically rational, so the integrality
statements can be tested exactly on random finite data. The script terminates with
\texttt{ALL CHECKS PASSED}.

\begin{lstlisting}[language=Python]
"""
Exact verification of the algebraic identities, integrality lemmas and
the construction used in the paper "An exponent window for Erdos Problem
#265: cross-tail integrality and controlled two-target greedy expansions".
Section/lemma numbers below refer to that paper.

All arithmetic is exact (fractions.Fraction).  For a FINITE strictly
increasing sequence the reciprocal sums are automatically rational, so
every integrality statement of Section 2 and Section 4 can be tested
exactly on random finite data.  (For finite data the tails eventually
vanish; positivity statements are only claimed while at least two terms
remain in the tail.)

Checks performed
  (C1) diagonal identity  sum_k 1/((a_k+h)(a_k+h')) = (R_h - R_h')/(h'-h)
  (C2) cross-tail identity E = R_h R_h' - (R_h-R_h')/(h'-h)
       equals the off-diagonal double sum
  (C3) integrality of q q' (h'-h) Pi_h Pi_h' E   and of
       q q' (h'-h) Pi_h Pi_h' * defect
  (C4) the bound E_n <= 2 R_h(n) R_h(n+1)            (h < h')
  (C5) the bound defect_m <= (p-m)/alpha_m^2 + R_h(p)/alpha_p  (m<=p)
  (C6) the multilinear Moebius identity for the fully off-diagonal sum
       of d >= 2 shifted tails, and the integrality of
       D * prod_i q_i Pi^(i)_n * O_n   (D = prod_{i<j}|h_j-h_i|)
  (C7) the bound O_n <= d! * G_n G_{n+1} ... G_{n+d-1}
  (C8) for Cantor's example a_n = C(n,2), n >= 3: integrality and
       positivity of N_n, M_n (Lemmas 2.3, 2.4) and the bound of Lemma 2.7(i)
  (C9) numerical facts used in the text: (1+2j)/2^j, Sylvester's
       telescoping identity, s_n^(1/2^n) -> E^2 = 1.59791..., the
       Badea comparison for negative integer shifts and the Duverney
       comparison of Example 6.1
  (C10) the Pell-type identity of Remark 7.1, as a rational identity
       and on the defect integers N_k of random finite sequences
  (C11) the identities and integrality for two non-proportional affine
       images U a_n + V, U' a_n + V' (Remark 3.5)
  (C12) the residual identity  Delta*E_n = (1-Delta R'_n)(R'_n/(1-Delta R'_n) - R_n),
       which identifies the cross tail with Kitamura's "second residual"
       for Delta = 1 (Remark 2.6)
  (C13) the construction of Theorem C: the one-step Lemma 5.1 (a)-(e) is
       checked in exact arithmetic along the first steps of the sequence
       built in the proof of Theorem C (beta = 5/4, K = 64), together with
       the bookkeeping inequalities used in that proof
"""
from fractions import Fraction as F
from itertools import permutations, product
from math import factorial, prod
import random

random.seed(20261028)

def tails(seq, h):
    """R[n] = sum_{k>=n} 1/(seq[k]+h), n = 0..len (R[len]=0)."""
    R = [F(0)] * (len(seq) + 1)
    for n in range(len(seq) - 1, -1, -1):
        R[n] = R[n + 1] + F(1, seq[n] + h)
    return R

def prefix_products(seq, h):
    P = [1] * (len(seq) + 1)
    for n in range(len(seq)):
        P[n + 1] = P[n] * (seq[n] + h)
    return P

def random_increasing(length, lo=3, growth=3.0):
    s, x = [], lo
    for _ in range(length):
        x = max(x + 1, int(x * random.uniform(1.05, growth)) + random.randint(0, 3))
        s.append(x)
    return s

# ---------------------------------------------------------------- d = 2
def check_two_shifts(trials=300):
    for _ in range(trials):
        L = random.randint(3, 9)
        seq = random_increasing(L)
        h = random.randint(-2, 3)
        hp = h + random.randint(1, 4)
        if min(seq) + h <= 0:
            continue
        Dlt = hp - h
        R, Rp = tails(seq, h), tails(seq, hp)
        P, Pp = prefix_products(seq, h), prefix_products(seq, hp)
        q, qp = R[0].denominator, Rp[0].denominator
        for n in range(L):
            # (C1)
            diag = sum(F(1, (seq[k] + h) * (seq[k] + hp)) for k in range(n, L))
            assert diag == (R[n] - Rp[n]) / Dlt
            # (C2)
            E = R[n] * Rp[n] - (R[n] - Rp[n]) / Dlt
            off = sum(F(1, (seq[j] + h) * (seq[k] + hp))
                      for j in range(n, L) for k in range(n, L) if j != k)
            assert E == off
            # (C3)
            for val in (E, diag):
                z = q * qp * Dlt * P[n] * Pp[n] * val
                assert z.denominator == 1, (seq, h, hp, n, z)
            if L - n >= 2:
                assert E > 0
            # (C4)
            if n + 1 <= L:
                assert E <= 2 * R[n] * R[n + 1]
            # (C5)
            for p in range(n, L + 1):
                bound = F(p - n, (seq[n] + h) ** 2)
                if p < L:
                    bound += R[p] / (seq[p] + h)
                assert diag <= bound
    print("d=2 identities (C1)-(C5) verified on", trials, "random finite sequences")

# ---------------------------------------------------------------- general d
def set_partitions(items):
    if not items:
        yield []
        return
    first, rest = items[0], items[1:]
    for part in set_partitions(rest):
        yield [[first]] + part
        for i in range(len(part)):
            yield part[:i] + [[first] + part[i]] + part[i + 1:]

def moebius(partition):
    return prod((-1) ** (len(B) - 1) * factorial(len(B) - 1) for B in partition)

def check_multilinear(trials=60):
    for _ in range(trials):
        d = random.randint(2, 4)
        L = random.randint(d, 7)
        seq = random_increasing(L)
        hs = sorted(random.sample(range(-2, 6), d))
        if min(seq) + hs[0] <= 0:
            continue
        Rs = [tails(seq, h) for h in hs]
        Ps = [prefix_products(seq, h) for h in hs]
        qs = [R[0].denominator for R in Rs]
        Dfac = prod(abs(hs[j] - hs[i]) for i in range(d) for j in range(i + 1, d))
        for n in range(L):
            idx = list(range(n, L))
            # fully off-diagonal sum, brute force
            O = F(0)
            for ks in product(idx, repeat=d):
                if len(set(ks)) == d:
                    O += prod(F(1, seq[k] + hs[i]) for i, k in enumerate(ks))
            # Moebius / partial-fraction expression
            expr = F(0)
            for part in set_partitions(list(range(d))):
                term = F(moebius(part))
                for B in part:
                    SB = F(0)
                    for i in B:
                        c = F(1)
                        for j in B:
                            if j != i:
                                c /= (hs[j] - hs[i])
                        SB += c * Rs[i][n]
                    term *= SB
                expr += term
            assert expr == O, (seq, hs, n)
            z = Dfac * prod(qs[i] * Ps[i][n] for i in range(d)) * O
            assert z.denominator == 1
            if L - n >= d:
                assert O > 0
            # (C7) with G = tail for the smallest shift
            G = Rs[0]
            bnd = factorial(d) * prod(G[min(n + j, L)] for j in range(d))
            assert O <= bnd
    print("multilinear identity, integrality and bound (C6)-(C7) verified on", trials, "random cases")

# ---------------------------------------------------------------- Cantor example
def check_cantor():
    from math import comb
    q, qp = 1, 9          # sum 1/C(n,2) = 1, sum 1/(C(n,2)-1) = 11/9  (n >= 3)
    R = lambda n: F(2, n - 1)
    Rp = lambda n: F(2, 3) * (F(1, n - 2) + F(1, n - 1) + F(1, n))
    assert R(3) == 1 and Rp(3) == F(11, 9)
    Pi, Pip = 1, 1
    for n in range(3, 60):
        a = comb(n, 2)
        E = R(n) * Rp(n) + R(n) - Rp(n)
        dl = Rp(n) - R(n)
        assert (q * qp * Pi * Pip * E).denominator == 1 and E > 0
        assert (q * qp * Pi * Pip * dl).denominator == 1 and dl > 0
        assert E <= 2 * Rp(n) * Rp(n + 1)
        Pi *= a
        Pip *= a - 1
    print("Cantor example a_n = C(n,2): integrality and positivity of N_n, M_n and the bound of Lemma 2.7(i) verified for n = 3..59")

# ---------------------------------------------------------------- numerics used in the text
def check_numerics():
    # coefficients (1+2j)/2^j used in Lemma 3.3
    assert [F(1 + 2 * j, 2 ** j) for j in range(1, 5)] == [F(3, 2), F(5, 4), F(7, 8), F(9, 16)]
    # Sylvester: s_0=2, s_{n+1}=s_n^2-s_n+1, sum 1/s_n = 1 (telescoping), s_n^{1/2^n} -> E^2
    from decimal import Decimal, getcontext
    getcontext().prec = 60
    s, tot = 2, F(0)
    for n in range(12):
        tot += F(1, s)
        # telescoping identity: 1 - sum_{k<=n} 1/s_k = 1/(s_{n+1}-1)
        nxt = s * s - s + 1
        assert 1 - tot == F(1, nxt - 1)
        if n >= 6:
            val = Decimal(s).ln() / Decimal(2 ** n)
            approx = val.exp()
        s = nxt
    E2 = approx
    assert abs(E2 - Decimal("1.5979102")) < Decimal("1e-6"), E2
    # Badea comparison for negative integer shifts t=-m: (s_{n+1}-m) - ((s_n-m)^2-(s_n-m)+1) = m(2 s_n - m - 2) > 0
    s = 2
    for n in range(8):
        nxt = s * s - s + 1
        for m in range(1, 5):
            if s - m >= 1:
                lhs = (nxt - m) - ((s - m) ** 2 - (s - m) + 1)
                assert lhs == m * (2 * s - m - 2)
        s = nxt
    # Duverney comparison (Example 6.1): u_n = w s_n + u satisfies
    # u_{n+1} - (u_n^2/w - u_n + w) = -u (2 s_n + u/w - 2), nonzero for u != 0 and large n
    s = 2
    for n in range(8):
        nxt = s * s - s + 1
        for w in range(1, 5):
            for u in range(-3, 4):
                un, un1 = w * s + u, w * nxt + u
                lhs = un1 - (F(un * un, w) - un + w)
                assert lhs == -u * (2 * s + F(u, w) - 2)
        s = nxt
    print("numerical constants verified: E^2 ~", str(E2)[:12])

# ---------------------------------------------------------------- Pell-type identity (Remark 7.1)
def check_pell(trials=2000):
    """N_{k+1}=y_k N_k - P_k recursion => y_{k+1}N_{k+1} = y_k^2 N_k - y_k N_{k+1} + N_{k+2};
    with z_k = 2 alpha_k + Delta, 4 y_k = z_k^2 - Delta^2, verify
    u_k^2 - 4 N_k N_{k+1} z_{k+1}^2 = c_k  (as an identity in rationals)."""
    for _ in range(trials):
        Dl = random.randint(1, 6)
        N0, N1, N2 = (random.randint(1, 50) for _ in range(3))
        alpha = random.randint(1, 10 ** 6)
        z0 = 2 * alpha + Dl
        y0 = F(z0 * z0 - Dl * Dl, 4)
        y1 = (y0 * y0 * N0 - y0 * N1 + N2) / F(N1)
        z1sq = 4 * y1 + Dl * Dl
        u = N0 * z0 * z0 - N0 * Dl * Dl - 2 * N1
        c = 4 * N1 * N1 - 4 * N0 * N1 * Dl * Dl - 16 * N0 * N2
        assert u * u - 4 * N0 * N1 * z1sq == c
    # and on genuine data: the defect integers of random finite sequences with h'-h = Delta
    for _ in range(200):
        L = random.randint(4, 8)
        seq = random_increasing(L)
        h = random.randint(-1, 2)
        Dl = random.randint(1, 3)
        if min(seq) + h <= 0:
            continue
        al = [x + h for x in seq]
        be = [x + Dl for x in al]
        y = [a * b for a, b in zip(al, be)]
        Dtail = [F(0)] * (L + 1)
        for n in range(L - 1, -1, -1):
            Dtail[n] = Dtail[n + 1] + F(1, y[n])
        S, Sp = sum(F(1, a) for a in al), sum(F(1, b) for b in be)
        Q = S.denominator * Sp.denominator
        PP = [1] * (L + 1)
        for n in range(L):
            PP[n + 1] = PP[n] * y[n]
        Nn = [Dl * Q * PP[n] * Dtail[n] for n in range(L + 1)]
        for n in range(L):
            assert Nn[n].denominator == 1
        for k in range(L - 2):
            z0, z1 = 2 * al[k] + Dl, 2 * al[k + 1] + Dl
            u = Nn[k] * z0 * z0 - Nn[k] * Dl * Dl - 2 * Nn[k + 1]
            c = 4 * Nn[k + 1] ** 2 - 4 * Nn[k] * Nn[k + 1] * Dl * Dl - 16 * Nn[k] * Nn[k + 2]
            assert u * u - 4 * Nn[k] * Nn[k + 1] * z1 * z1 == c
    print("Pell-type identity of Remark 7.1 verified (symbolic-random and on genuine defect data)")

# ---------------------------------------------------------------- affine images (Remark 3.5)
def check_affine(trials=300):
    done = 0
    while done < trials:
        L = random.randint(3, 8)
        seq = random_increasing(L)
        U, Up = random.randint(1, 4), random.randint(1, 4)
        V, Vp = random.randint(-2, 5), random.randint(-2, 5)
        Dst = Up * V - U * Vp
        if Dst == 0:
            continue
        al = [U * x + V for x in seq]
        be = [Up * x + Vp for x in seq]
        if min(al) <= 0 or min(be) <= 0:
            continue
        done += 1
        R, Rp = [F(0)] * (L + 1), [F(0)] * (L + 1)
        for n in range(L - 1, -1, -1):
            R[n] = R[n + 1] + F(1, al[n])
            Rp[n] = Rp[n + 1] + F(1, be[n])
        Q = R[0].denominator * Rp[0].denominator
        P, Pp = 1, 1
        for n in range(L):
            diag = sum(F(1, al[k] * be[k]) for k in range(n, L))
            assert diag == (Up * Rp[n] - U * R[n]) / Dst
            E = R[n] * Rp[n] - diag
            off = sum(F(1, al[j] * be[k]) for j in range(n, L) for k in range(n, L) if j != k)
            assert E == off
            assert (abs(Dst) * Q * Q * P * Pp * E).denominator == 1
            assert (abs(Dst) * Q * P * Pp * diag).denominator == 1
            P *= al[n]
            Pp *= be[n]
    print("affine-image identities and integrality (Remark 3.5) verified on", trials, "cases")

# ---------------------------------------------------------------- residual form (Remark 2.6)
def check_residual(trials=300):
    done = 0
    while done < trials:
        L = random.randint(3, 8)
        seq = random_increasing(L, lo=5)
        h = random.randint(-2, 2)
        Dl = random.randint(1, 4)
        al = [x + h for x in seq]
        if min(al) <= 0:
            continue
        be = [x + Dl for x in al]
        R, Rp = [F(0)] * (L + 1), [F(0)] * (L + 1)
        for n in range(L - 1, -1, -1):
            R[n] = R[n + 1] + F(1, al[n])
            Rp[n] = Rp[n + 1] + F(1, be[n])
        for n in range(L):
            if 1 - Dl * Rp[n] == 0:
                continue
            E = R[n] * Rp[n] - (R[n] - Rp[n]) / Dl
            rho = Rp[n] / (1 - Dl * Rp[n]) - R[n]
            assert Dl * E == (1 - Dl * Rp[n]) * rho
        done += 1
    print("residual identity (C12) verified on", trials, "random cases")

# ---------------------------------------------------------------- construction (Theorem C, Lemma 5.1)
def check_construction(steps=11):
    from decimal import Decimal, getcontext
    import math
    getcontext().prec = 80
    K = 64
    beta = F(5, 4)
    N = 10 ** 11

    def dec(q):
        return Decimal(q.numerator) / Decimal(q.denominator)

    def lg(q):            # natural log of a positive Fraction
        return math.log(q.numerator) - math.log(q.denominator)

    def after(T, D, a):
        Tn = T - F(1, a)
        Dn = D - F(1, a * (a - 1))
        Gn = Tn * Tn - Dn
        return Tn, Dn, Gn

    T = F(1, N)
    s1 = F(Decimal(N) ** (-Decimal(beta.numerator) / Decimal(beta.denominator)))   # s_1 ~ N^{-beta}
    D = T * T - 2 * T * s1
    T1, D1 = T, D
    seq = []
    lam_prev = None
    for n in range(1, steps + 1):
        G = T * T - D
        s = G / (2 * T)
        # invariant (8) and T <= K^-2
        assert 0 < T <= F(1, K * K) and K * T * T <= s <= T / K
        lam, sig = -lg(T), -lg(s)
        assert abs(sig - float(beta) * lam) <= 1.0                           # |nu_n - beta*lambda_n| <= 1 (floating point)
        # target y: rational approximation of s^beta (either side); Lemma 5.1 allows any real y with (9)
        y = F(dec(s) ** (Decimal(beta.numerator) / Decimal(beta.denominator)))   # y = 80-digit rational approximation of s^beta
        assert 8 * T ** 3 / s <= y <= s / (4 * K)                            # condition (9), exact
        # Lemma 5.1: find a > 1/T with s' in [y, 2y]; s'(a) is decreasing on (1/T, 2/T]
        lo = T.denominator // T.numerator + 1
        hi = (2 * T.denominator) // T.numerator
        def sp(a):
            Tn, Dn, Gn = after(T, D, a)
            return Gn / (2 * Tn)
        assert sp(lo) >= y and sp(hi) < y
        while hi - lo > 1:
            mid = (lo + hi) // 2
            if sp(mid) >= y:
                lo = mid
            else:
                hi = mid
        a = lo
        Tn, Dn, Gn = after(T, D, a)
        spn = Gn / (2 * Tn)
        # conclusions (a)-(e) of Lemma 5.1
        assert Tn > 0 and Dn > 0 and Gn > 0                                  # (a)
        assert (1 - F(4, K)) * s <= Tn <= (1 + F(6, K)) * s                  # (b)
        assert y <= spn <= 2 * y                                              # (c)
        assert Tn <= F(1, K * K) and K * Tn * Tn <= spn <= Tn / K            # (d)
        assert F(1) / T < a <= (1 + F(4, K)) / T                              # (e)
        if seq:
            assert a > seq[-1]
        seq.append(a)
        T, D = Tn, Dn
    # bookkeeping: the partial sums reproduce the rational targets exactly
    assert T1 - sum(F(1, a) for a in seq) == T
    assert D1 - sum(F(1, a * (a - 1)) for a in seq) == D
    ratios = [math.log(seq[i + 1]) / math.log(seq[i]) for i in range(len(seq) - 1)]
    norm = [math.log(seq[i]) / float(beta) ** (i + 1) for i in range(len(seq))]
    assert all(abs(r - float(beta)) < 0.05 for r in ratios[2:])
    print("construction of Theorem C (beta=" + str(beta) + ", K=64): Lemma 5.1 (a)-(e) verified for", steps,
          "steps; log a_n / beta^n =", ", ".join(f"{v:.3f}" for v in norm[-4:]))

if __name__ == "__main__":
    check_two_shifts()
    check_multilinear()
    check_cantor()
    check_numerics()
    check_pell()
    check_affine()
    check_residual()
    check_construction()
    print("ALL CHECKS PASSED")
\end{lstlisting}

\begin{thebibliography}{99}

\bibitem{Badea87}
C.~Badea, \emph{The irrationality of certain infinite series},
Glasgow Math.\ J.\ \textbf{29} (1987), no.~2, 221--228.

\bibitem{Bloom265}
T.~F.~Bloom, \emph{Erd\H{o}s Problem \#265},
\url{https://www.erdosproblems.com/265} (accessed October 28, 2026).

\bibitem{Duverney01}
D.~Duverney, \emph{Irrationality of fast converging series of rational numbers},
J.\ Math.\ Sci.\ Univ.\ Tokyo \textbf{8} (2001), 275--316.

\bibitem{EPAD265}
Erd\H{o}s Problem a Day, \emph{Erd\H{o}s \#265 -- working report} (AI-generated, 2026),
\url{https://erdosproblemaday.com/report/265}.

\bibitem{Er85}
P.~Erd\H{o}s, \emph{E.~Straus (1921--1983)}, Rocky Mountain J.\ Math.\ \textbf{15}
(1985), no.~2, 331--341.

\bibitem{Er88}
P.~Erd\H{o}s, \emph{On the irrationality of certain series: problems and results},
in: New Advances in Transcendence Theory (Durham, 1986), A.~Baker (ed.),
Cambridge Univ.\ Press, 1988, pp.~102--109.

\bibitem{ErGr80}
P.~Erd\H{o}s and R.~L.~Graham, \emph{Old and New Problems and Results in Combinatorial
Number Theory}, Monographies de L'Enseignement Math\'ematique \textbf{28},
Geneva, 1980.

\bibitem{ErSt64}
P.~Erd\H{o}s and E.~G.~Straus, \emph{On the irrationality of certain Ahmes series},
J.\ Indian Math.\ Soc.\ (N.S.) \textbf{27} (1964), 129--133.

\bibitem{FC6839}
google-deepmind/formal-conjectures, \emph{Pull request \#6839: add Erd\H{o}s 265 with
external proof reference} (October 2026),
\url{https://github.com/google-deepmind/formal-conjectures/pull/6839}.

\bibitem{Kit26}
K.~Kitamura, \emph{erdos-265-lean: a Lean~4 proof that the two rational reciprocal sums
in Erd\H{o}s Problem~265 force $a_n^{1/2^n}\to1$}, GitHub repository, commit
\texttt{16db9f29} (September 2026),
\url{https://github.com/KitaKen1/erdos-265-lean}.

\bibitem{Koizumi25}
J.~Koizumi, \emph{Irrationality of the reciprocal sum of doubly exponential sequences},
arXiv:2504.05933 (2025).

\bibitem{Kov24}
V.~Kova\v{c}, \emph{On simultaneous rationality of two Ahmes series},
arXiv:2406.17593v1 (2024).

\bibitem{KT25}
V.~Kova\v{c} and T.~Tao, \emph{On several irrationality problems for Ahmes series},
Acta Math.\ Hungar.\ (2025), doi:10.1007/s10474-025-01528-0; arXiv:2406.17593v2--v4.

\bibitem{Sylvester1880}
J.~J.~Sylvester, \emph{On a point in the theory of vulgar fractions},
Amer.\ J.\ Math.\ \textbf{3} (1880), 332--335.

\bibitem{Vardi91}
I.~Vardi, \emph{Computational Recreations in Mathematica}, Addison--Wesley, 1991.

\end{thebibliography}

\end{document}
